\part{Characteristic Classes}

\chapter{Stiefel-Whitney Classes}

\section{Axiomatic definitions}

	\begin{defn} \label{SW_classes}
		For a topological space $B$ and a $n$-dimensional vector bundle $B(\xi)=p:E\to B$ ($B(\xi)$ will denote the base space of the bundle later when we want to emphasize the base space of the bundle), we define the \emph{Stiefel-Whitney classes} $w_i(\xi)$ of the bundle to be a collection of modulo 2 cohomology classes having the following properties: 
		\begin{enumerate}
			\item $w_i(\xi)\in H^i(B(\xi);\mathbb{Z}/2\mathbb{Z})$, in particular $w_0(\xi)=1\in H^0(B(\xi);\mathbb{Z}/2\mathbb{Z})$, and $w_i(\xi)=0$ for $i>n$. 
			\item The Stiefel-Whitney classes are natural. That is, for a bundle map $f:B(\xi)\to B(\eta)$, we have $w_i(\xi)=f^* w_i(\eta)$. 
			\item The Stiefel-Whitney classes of the direct sum bundle can be computed as follows: 
				\[ w_k(\xi\oplus\eta)=\sum_{i=0}^k w_i(\xi)\smile w_{k-i}(\eta) \]
			\item For the canonical line bundle over $\mathbb{RP}^1$, $\gamma_1^1$, we have $w_1(\gamma_1^1)$ nontrivial. 
		\end{enumerate}
	\end{defn}

	\begin{remark}
		Even though these definitions works well for general topological spaces, we shall temporarily use them only on manifolds, where the notion of vector bundles works best. As the most natural bundle of a manifold $M$, the tangent bundle $\tau_M$, the Stiefel-Whitney class of $\tau_M$ is often denoted by $w(M)$. 
	\end{remark}

	\begin{defn}
		We define $H^\Pi(X;\mathbb{Z}/2\mathbb{Z})$ to be the ring consisting of all formal infinite series 
		\[ a=a_0+a_1+\dots \]
		where $a_i\in H^i(X;\mathbb{Z}/2\mathbb{Z})$, and the obvious product defined on it. Then we define the total Stiefel-Whitney class to be $w(\xi)=w_0(\xi)+w_1(\xi)+\dots$. \\
		Under this definition, the third axiom could be written simply as $w(\xi\oplus\eta)=w(\xi)w(\eta)$. 
	\end{defn}

	\begin{remark}
		This ring has properties similar to the formal series ring. For example the units of this ring are $a=1+a_1+\dots$. 
	\end{remark}

	We have the following immediate consequences of the axioms: 
	
	\begin{coro}
		\begin{enumerate}
			\item If $\xi\simeq\eta$, then $w_i(\xi)=w_i(\eta)$ under the natural equivalence of the cohomology groups. 
			\item If $\varepsilon$ is trivial, then $w_i(\varepsilon)=0$ for $i>0$. 
			\item If $\varepsilon$ is trivial, then for any vector bundle $\xi$, $w_i(\xi\oplus\varepsilon)=w_i(\xi)$. 
			\item If $\xi$ is an $n$-dimensional bundle equipped with an Euclidean metric having a non-vanishing section, then $w_n(\xi)=0$. 
			\item $w(\xi)$ are units in $H^\Pi(B(\xi);\mathbb{Z}/2\mathbb{Z})$. 
		\end{enumerate}
	\end{coro}
	\begin{proof}
		Only the fourth property is non-trivial. To see this, let $\varepsilon$ be the trivial bundle generated by the non-vanishing section, then notice $\xi$ splits as $\varepsilon\oplus\varepsilon^\bot$. 
	\end{proof}

	\begin{defn}
		Being a unit, the inverse of $w(\xi)$ in $H^\Pi(B(\xi);\mathbb{Z}/2\mathbb{Z})$ is defined to be $\bar{w}(\xi)$. 
	\end{defn}

	\begin{coro}
		If we have $\xi\oplus\eta$ a trivial bundle on $B$, then $w(\xi)=\bar{w}(\eta)$. 
	\end{coro}
	\begin{remark}
		Often, the above corollary is used when $\xi$ is the tangent bundle and $\eta$ the normal bundle. 
	\end{remark}

\section{Elementary computations in $S^n$ and $\mathbb{RP}^n$}

	\begin{coro}
		$w(\tau_{S^n})=1$. 
	\end{coro}
	\begin{proof}
		Consider the embedding $S^n\inj\mathbb{R}^{n+1}$. Then we have the normal bundle $\nu$ and the direct sum $\nu\oplus\tau$ both trivial. 
	\end{proof}

	\begin{prop}
		We have $H^*(\mathbb{RP}^n;\mathbb{Z}/2\mathbb{Z})=\mathbb{Z}/2\mathbb{Z}[a]$, where the degree of $x$ is one. 
	\end{prop}
	\begin{proof}
		We use the standard cell structure on $\mathbb{RP}^n$. Then, since the mapping degree of the characteristic map of the cells are $2$, we have the map $C^i\to C^{i+1}$ the trivial map. This implies $H^i(\mathbb{RP}^n;\mathbb{Z}/2\mathbb{Z})=\mathbb{Z}/2\mathbb{Z}$ for $0\le i\le n$. The cup product structure follows directly. 
	\end{proof}

	\begin{prop}
		The Stiefel-Whitney class of the canonical line bundle $\gamma^1_n$ of $\mathbb{RP}^n$ is given by $w(\gamma^1_n)=1+a$. 
	\end{prop}
	\begin{proof}
		We consider the standard inclusion $j:\mathbb{RP}^1\to\mathbb{RP}^n$, covered by the bundle inclusion $\gamma^1_1\to\gamma^1_n$. Thus $j^*w_1(\gamma_n^1)=w_1(\gamma_1^1)\ne0$, and since the bundle is $1$-dimensional, the higher Stiefel-Whitney classes vanish. 
	\end{proof}

	\begin{coro}
		Recall that the trivial line bundle of $\mathbb{RP}^n$ embeds into the $n+1$-dimensional trivial bundle $\gamma_n^1\inj\varepsilon_{n+1}$. Thus we have the Stiefel-Whitney classes of the orthogonal complement of $w(\gamma^\bot)=1+a+a^2+\dots+a^n$. 
	\end{coro}

	We finish this section with the computations of the Stiefel-Whitney classes of the tangent bundle of $\mathbb{RP}^n$. 

	\begin{lemma}
		We have the isomorphism $\tau_{\mathbb{RP}^n}\simeq\Hom(\gamma_n^1,\gamma^\bot)$. 
	\end{lemma}
	\begin{proof}
		A direct check using the standard construction. 
	\end{proof}

	\begin{lemma}
		We have $\tau\oplus\varepsilon^1\simeq(\gamma_n^1)^{\oplus (n+1)}$. 
	\end{lemma}
	\begin{proof}
		By the preceding lemma, we have
		\begin{align*}
			&\tau\oplus\varepsilon^1\\
			\simeq&\Hom(\gamma_n^1,\gamma^\bot)\oplus\Hom(\gamma_n^1,\gamma_n^1)\\
			\simeq&\Hom(\gamma_n^1,\gamma^\bot\oplus\gamma_n^1)\\
			\simeq&\Hom(\gamma_n^1,\varepsilon^{n+1})\\
			\simeq&\Hom(\gamma_n^1,\varepsilon)^{\oplus (n+1)}\\
			\simeq&(\gamma_n^1)^{\oplus (n+1)}
		\end{align*}
	\end{proof}

	\begin{coro}
		\[ w(\mathbb{RP}^n)=(1+a)^{n+1} \]
	\end{coro}
	\begin{proof}
		Directly compute using the preceding lemma. 
	\end{proof}

\section{The Thom isomorphism}

	\begin{defn}
		For a vector bundle $\xi=\pi:E\to B$, we denote $E_0$ the subspace of $E$ consisting of all the non-zero elements within each fibre. If we let $F$ be the fibre, then $F_0$ is similarly defined as $F-\{0\}$. \\
		We say this bundle is $R$-oriented if there is a family of $R$-coefficient cohomology class $u_x\in H^n(F_x,F_0;R)$ which is a generator and agrees with each other in the locally trivializable condition. \\
		If there is a class $u\in H^n(E,E_0;R)$ such that $u|_x=u_x$, we call $u$ the fundamental class of the bundle. \\
		Equivalently, we may consider the unit ball bundle and unit sphere bundle associated to $\xi$, say $B(\xi)$ and $S(\xi)$. Then since $(E,E_0)$ is homotopy equivalent to $(B(\xi),S(\xi))$, we may equivalently define the fundamental class in $H^n(B(\xi),S(\xi);R)$. 
	\end{defn}

	The main purpose of this section is to prove the Thom isomorphism theorem, and use it to derive some important properties: 
	\begin{thm} \label{Thom_isom}
		For an oriented $R$-bundle $\xi=p:E\to B$ having fundamental class $u\in H^n(E,E_0,R)$, we have an isomorphism $\smile u:H^i(E;R)\to H^{n+i}(E,E_0;R)$. We define $\varphi:H^i(B;R)\to H^{n+i}(E,E_0;R)$ be the canonical isomorphism composed with the isomorphism $p^*:H^i(B;R)\to H^i(E;R)$. In particular, since $\mathbb{Z}/2\mathbb{Z}$-bundles are always oriented, the above isomorphism is always available. We may also substitute all the above $H^*(E,E_0;R)$ into $H^*(B(\xi),S(\xi);R)$. 
	\end{thm}
	\begin{remark}
		For a vector bundle $\xi=p:E\to B$, we may define its Thom space $\Th(\xi)$ to be the one-point compactification of $E$. Then, notice that $\tilde{H}^*(\Th(\xi))\simeq H^*(E,E_0)$. Therefore the Thom isomorphism can also be expressed by $H^*(B)\simeq H^{*+n}(\Th(\xi))$. 
	\end{remark}

	We first recall the Leray-Hirsch theorem. 
	\begin{thm} \label{Leray-Hirsch}
		Suppose we have a fibration $F\inj E\surj B$ and its subfibration $F'\inj E'\surj B$, such that $H^i(F,F';R)$ is a finitely generated free $R$-module for each $i$. Assume further that there is are classes $c_j\in H^*(E,E';R)$ whose restrictions form a basis for $H^*(F,F';R)$, then $H^*(E,E';R)$ is a free $H^*(B;R)$ module with basis $\{c_j\}$, where $bc:=p^*(b)\smile c$. 
	\end{thm}
	
	We use this to prove the second part of the Thom isomorphism theorem directly. 
	\begin{proof}
		Notice $H^i(F,F';R)=R$ when $i=0,n$ and is $0$ otherwise. Thus $u\in H^*(E,E';R)$ is a basis for $H^*(F,F';R)$, and thus is a $H^*(B;R)$ basis of $H^*(E,E';R)$. 
	\end{proof}

	Then, we prove following property: 
	\begin{prop} \label{partial_fund_class}
		If $\xi$ is an $R$-oriented bundle, suppose either $R$ is a field or $B$ is compact, then the fundamental class always exists. 
	\end{prop}
	\begin{proof}
		We prove there is a unique $u\in H^n(E,E_0;R)$ that restricts to $u_x$. We finish the proof in four steps. \\
		The first step: when $\xi$ is a trivial bundle, this is obvious. \\
		The second step: when $B=B_1\cup B_2$, and $\xi$ is a trivial bundle when restricted to both $B_1$ and $B_2$, we let $E^1=p^{-1}(B_1)$ and $E^2=p^{-1}(B_2)$. Then, we have the exact sequence obtained by the Mayer-Vietoris sequence below: 
		\[ H^n(E,E_0;R)\lto H^n(E^1,E_0^1;R)\oplus H^n(E^2,E_0^2;R)\lto H^n(E^1\cap E^2,E_0^1\cap E_0^2;R) \]
		Now that we have a class $u_1$ and $u_2$ in $H^n(E^1,E_0^1;R)$ and $H^n(E^2,E_0^2;R)$ respectively, their restriction to $H^n(E^1\cap E^2,E_0^1\cap E_0^2;R)$ coincides by the uniqueness condition on $\xi|_{B_1\cap B_2}$ (note this is a trivial bundle). Thus by exactness, there is a class $u\in H^n(E,E_0;R)$ that restricts to $u_1$ and $u_2$, and thus restricts to all $u_x$. \\
		The third step: when $B$ is compact, we cover $B$ by open subsets $B_i$ such that $\xi$ is trivializable over $B_i$. Then by compactness, we can find a finite subcover of $B_i$. Then by the second step, we have the uniqueness and existence. \\
		Note here that we have already proven the proposition on compact spaces. 
		The fourth step: when $B$ is arbitrary, we attempt to prove $H^n(B;R)=\plim_{C\subseteq B} H^n(C;R)$, where $C$ is compact in $B$ and ordered by inclusion. To prove this, first notice that we have $\ilim H_n(C;R)\to H_n(B;R)$ is an isomorphism since homology is naturally compactly supported. Then, for field coefficients, $H^n(X;R)$ is naturally isomorphic to $\Hom(H_n(X;R);R)$ by the universal coefficient theorem applied to $R$-coefficient homology and cohomology, where the $\mathrm{Ext}(H_{n-1}(X;R),R)$ term vanishes since $H_{n-1}(X;R)$ is a free $R$-module. Thus 
		\begin{align*}
			&H^n(B;R) \\
			=&\Hom(H_n(B;R),R) \\
			\simeq&\Hom(\ilim H_n(C;R),R) \\
			\simeq&\plim\Hom(H_n(C;R),R) \\
			=&\plim H^n(C;R) \\
		\end{align*}
		Similarly $H^n(E,E_0;R)=\plim H^n(p^{-1}(C),p^{-1}(C)_0;R)$ since obviously every compact subspace of $(E,E_0)$ is contained in some $(p^{-1}(C),p^{-1}(C)_0)$. Then the classes in $H^n(p^{-1}(C),p^{-1}(C)_0;R)$ obviously agrees with inclusions by the uniqueness in compact case, thus put together to a class $u\in H^n(E,E_0;R)$. 
	\end{proof}

	With another technical lemma, we prove the case for arbitrary coefficients and base spaces. 
	\begin{lemma}
		Let $K$ and $K'$ be free complexes over $\mathbb{Z}$, and $f:K\to K'$ a chain map of degree $d$. Then if $f:H^*(K',F)\to H^*(K,F)$ is an isomorphism for every field $K$, then $f$ induces isomorphisms of homology and cohomology with arbitrary coefficients. 
	\end{lemma}
	\begin{proof}
		Without loss of generality, we suppose $d=0$. Recall the construction of the mapping cone $K^f$. It then fits in an exact sequence of complexes $0\lto K'\lto K^f\lto K\lto 0$. We prove that $H_n(K^f)=0$ for all $n$. Firstly, since $f:H^*(K';F)\to H^*(K;F)$ is an isomorphism, we have $f^*:\Hom(H_*(K'\otimes F),F)\to\Hom(H_*(K\otimes F),F)$ an isomorphism, thus $\Hom(H_*(K^f\otimes F),F)=0$ by the long exact sequence. Therefore $H_n(K^f\otimes F)=0$. In particular, $H_n(K^f\otimes\mathbb{Q})=0$, and then it suffices to prove the torsion group $H_n(K^f)=0$. This follows from considering $H_n(K^f\otimes\mathbb{F}_p)=0$ to eliminate $p$-torsion parts. 
	\end{proof}

	\begin{prop}
		For a bundle $\xi$, if $u$ is the fundamental class of $\xi$, then the correspondence $u\frown:H_{n+i}(E,E_0;\mathbb{Z})\to H_i(E;\mathbb{Z})$ is an isomorphism. 
	\end{prop}
	\begin{proof}
		We pick a cochain class $v$ representing $u$. Note $v\frown:C_{n+i}(E,E_0;\mathbb{Z})\to C_i(E;\mathbb{Z})$ is a chain map. Then, using the identity $\langle c,v\frown d\rangle=\langle c\smile v,d\rangle$ (here $c$ is a cochain class and $d$ a chain class), we have the cohomology mapping induced by $v\frown$ is $\smile u$. Then the conclusion follows from the preceding lemma. 
	\end{proof}

	We are finally ready to prove the full version of \cref{partial_fund_class}
	\begin{thm} \label{fund_class}
		If $\xi$ is an $R$-oriented bundle, then the fundamental class always exists. 
	\end{thm}
	\begin{proof}
		The first three steps are the same as the previous partial proof. For the fourth step, we have $H_{n-1}(E,E_0;\mathbb{Z})\simeq H_{-1}(E;\mathbb{Z})=0$, and thus $H^n(E,E_0;R)=\Hom(H_n(E,E_0;\mathbb{Z}),R)$. The rest is also the same. 
	\end{proof}

	\begin{coro}
		We have a fundamental class in $u\in H^n(E,E_0;\mathbb{Z}/2\mathbb{Z})$ inducing the Thom isomorphism $\smile u:H^i(E;\mathbb{Z}/2\mathbb{Z})\to H^i(E,E_0;\mathbb{Z}/2\mathbb{Z})$. 
	\end{coro}
	\begin{proof}
		A bundle is always $\mathbb{Z}/2\mathbb{Z}$ oriented. 
	\end{proof}

\section{Existence and uniqueness of the Stiefel-Whitney classes}
	
	Before we construct the Stiefel-Whitney classes explicitly, we first review some properties of the Steenrod operations of cohomology. 
	\begin{thm} \label{Steenrod_op}
		There exists homomorphisms $Sq^i:H^n(X,Y;\mathbb{Z}/2\mathbb{Z})\to H^{n+i}(X,Y;\mathbb{Z}/2\mathbb{Z})$ which can be defined using the following axioms: 
		\begin{enumerate}
			\item $Sq^0=\id$. 
			\item $Sq^i$ is natural. 
			\item $Sq^i(\alpha\smile\beta)=\sum_{j=0}^i Sq^j(\alpha)\smile Sq^{i-j}(\beta)$. 
			\item $\Sigma Sq^i=Sq^i\Sigma$, where $\Sigma$ is the natural suspension isomorphism \\
			$H^n(X,Y;\mathbb{Z}/2\mathbb{Z})\to H^{n+1}(\Sigma X,\Sigma Y;\mathbb{Z}/2\mathbb{Z})$. 
			\item $Sq^i(\alpha)=\alpha^2$ if $i=n$ for $\alpha\in H^n(X,Y;\mathbb{Z}/2\mathbb{Z})$ and is $0$ for $i>n$. 
			\item $Sq^1$ coincides with the Bockstein homomorphism associated to the sequence
			\[ 0\lto \mathbb{Z}/2\mathbb{Z}\lto \mathbb{Z}/4\mathbb{Z}\lto \mathbb{Z}/2\mathbb{Z}\lto 0 \]
			\item $Sq^i$ satisfies the Adem relations
			\[ Sq^a Sq^b=\sum_j\binom{b-j-1}{a-2j}Sq^{a+b-j}Sq^j,~a<2b \]
		\end{enumerate}
		Moreover, we define the total steenrod square 
		\[ Sq(a)=Sq^0(a)+\dots+Sq^n(a) \] 
		for $a\in H^n(X,Y;\mathbb{Z}/2\mathbb{Z})$. 
	\end{thm}

	\begin{const} \label{Cons_SW}
		For a vector bundle $\xi=\pi:E\to B$, we let $w_i(\xi)=\varphi^{-1}Sq^i\varphi(1)=\varphi^{-1}Sq^i(u)\in H^i(B;\mathbb{Z}/2\mathbb{Z})$, where $\varphi:H^i(B;\mathbb{Z}/2\mathbb{Z})\to H^{i+n}(E,E_0;\mathbb{Z}/2\mathbb{Z})$ is the Thom isomorphism (\cref{Thom_isom}), $Sq^i$ the Steenrod square operation, and $u$ the fundamental class. Or, equivalently, we define $w(\xi)=\varphi^{-1}Sq\varphi(1)=\varphi^{-1}Sq(u)$. \\
		We verify the axioms. 
		\begin{enumerate}
			\item Obviously $w_i(\xi)\in H^i(B;\mathbb{Z}/2\mathbb{Z})$, and $w_0=1$, $w_i=0$ for $i>n$. 
			\item Since both the Thom isomorphism and the Steenrod squares are natural. 
			\item Recall that for $\xi=p:E\to B$ and $\xi'=p':E'\to B$, we have the definition of $\xi\oplus\xi'$ the pullback bundle of $\xi\times\xi'=p\times p':E\times E'\to B\times B$ along $\Delta:B\to B\times B$. So to verify the third axiom, we first calculate the Stiefel-Whitney class of $\xi\times\xi'$. \\
			Consider the fundamental classes $u\in H^m(E,E_0;\mathbb{Z}/2\mathbb{Z})$ and $u'\in H^n(E',E'_0;\mathbb{Z}/2\mathbb{Z})$. Thus we have the cross product $u\times u'=\pi_1^*(u)\smile \pi_2^*(u')\in H^{m+n}(E\times E',E\times E'_0\cup E_0\times E';\mathbb{Z}/2\mathbb{Z})$. Notice then that $E\times E'_0\cup E_0\times E'$ are exactly the non-zero vectors in $E\times E'$, and we claim that $u\times u'$ is the fundamental class of $\xi\times\xi'$. This is verified directly by the naturality of cross products. \\
			Then, by the construction of the Thom isomorphism, we have $\varphi(\alpha\times\alpha')=\varphi(\alpha)\times\varphi(\alpha')$, where $\varphi$ denotes the three distinct Thom isomorphisms (note that there are no issues with signs since we are working modulo 2). \\
			Thus finally, we have
			\[ w(\xi\times\xi')=\varphi^{-1}(Sq(u\times u'))=\varphi^{-1}(Sq(u))\times\varphi^{-1}(Sq(u'))=w(\xi)\times w(\xi') \]
			Pulling back along the diagonal gives $w(\xi\oplus\xi')=w(\xi)\smile w(\xi')$. 
			\item Let $\gamma_1^1$ be the twisted line bundle over $\mathbb{RP}^1$ as usual, then the vectors with length $\le 1$ in $E=E(\gamma_1^1)$ is obviously a mobius band $M$. We suppose it is bound by the circle $\partial M$. Since $M$ is a deformation retract of $E$, and $\partial M$ a deformation retract of $E_0$, we have $H^*(M,\partial M;\mathbb{Z}/2\mathbb{Z})\simeq H^*(E,E_0;\mathbb{Z}/2\mathbb{Z})$. On the other hand, by definition $H^*(\mathbb{RP}^2,\mathbb{D}^2;\mathbb{Z}/2\mathbb{Z}/2\mathbb{Z})\simeq H^*(M,\partial M)$. Since $\mathbb{D}^2$ is contractible, we have \\
			$H^i(E,E_0;\mathbb{Z}/2\mathbb{Z})\simeq H^i(\mathbb{RP}^2;\mathbb{Z}/2\mathbb{Z})$, and thus the fundamental class $u$ must corresponds to the only non-zero generator $a\in H^1(\mathbb{RP}^2)$. Thus $Sq^1(u)=u\smile u$ corresponds to $a\smile a$. But since $a\smile a\ne0$, $w_1(\gamma_1^1)\ne0$. 
		\end{enumerate}
	\end{const}

	Now we prove the uniqueness. But before this, we need to cite an important theorem. 
	\begin{defn}\label{Defn_Grassmannian}
		For positive integers $n$ and $k$ and a field $\mathbb{F}$(usually $\mathbb{F}$ is equal to $\mathbb{R}$ or $\mathbb{C}$), we define the space $Gr_n(\mathbb{F}^{n+k})$ to be the quotient space $\Hom_\mathbb{F}(\mathbb{F}^n,\mathbb{F}^{n+k})/R$, where $R$ is the equivalence relation where $f=g$ iff $\img f=\img g$. \\
		Then we define $Gr_n(\mathbb{F}):=Gr_n(\mathbb{F}^\infty)=\ilim_kGr_n(\mathbb{F}^{n+k})$. Notice that we have a $n$-bundle on $Gr_n(\mathbb{F}^{n+k})$ constructed as follows. Let $X=\{(f,x)\in Gr_n(\mathbb{F}^{n+k})\times\mathbb{F}^{n+k}|x\in\img f\}$, then $(f,x)\mapsto f$ is an $n$-bundle on $Gr_n(\mathbb{F}^{n+k})$. This put together to be an $n$-bundle on $Gr_n(\mathbb{F})$, which we will denote by $\eta^n$. 
	\end{defn}
	\begin{thm}\label{Real_Grassmannian}
		For a paracompact topological space $X$, an $n$-dimensional vector bundle $p:E\to X$ corresponds to a homotopy class of maps $f:X\to Gr_n(\mathbb{R})$. Moreover, this bundle is the pullback bundle of the canonical $n$-bundle $\eta^n$ over $Gr_n(\mathbb{R})$. \\
		Or equivalently, we have a homotopy equivalence $Gr_n(\mathbb{R})\simeq BO(n)$. 
	\end{thm} 

	Then, we have an important observation that will allow us to finish the proof. 
	\begin{thm}\label{SW_gen}
		We have the cohomology $H^*(Gr_n(\mathbb{R});\mathbb{Z}/2\mathbb{Z})$ freely generated as a polynomial over $\mathbb{Z}/2\mathbb{Z}$ by the Stiefel-Whitney classes of the canonical $n$-bundle $\eta^n$. 
	\end{thm}
	\begin{proof}
		We first prove the Stiefel-Whitney classes $w_i(\eta^n)$ are algebraically independent. We consider the classifying map of the $n$-fold product $\xi=\gamma^1\times\cdots\times\gamma^1=\pi_1^{-1}(\gamma^1)\oplus\cdots\oplus\pi_n^{-1}(\gamma^1)$ on the space $\mathbb{RP}^\infty\times\cdots\times\mathbb{RP}^\infty$. Now, the classifying map $g:\mathbb{RP}^\infty\times\cdots\times\mathbb{RP}^\infty\to Gr_n(\mathbb{R})$ gives $w_i(\xi)=g^*w_i(\eta^n)$. Recall $H^*(\mathbb{RP}^\infty;\mathbb{Z}/2\mathbb{Z})=\mathbb{Z}/2\mathbb{Z}[a]$ on one generator having degree 1, thus $H^*(\mathbb{RP}^\infty\times\cdots\times\mathbb{RP}^\infty;\mathbb{Z}/2\mathbb{Z})=\mathbb{Z}/2\mathbb{Z}[a_1,\dots,a_n]$, all having degree 1. But note that $w(\xi)=w(\gamma^1)^n=(1+a_1)\cdots(1+a_n)$, thus $w_i(\xi)$ is the $i$-th symmetric polynomial. But it is well-known that they are algebraically independent, and thus $w_i(\eta^n)$ must be algebraically independent. \\
		Then, we prove that they indeed generate the whole cohomology group of $H^*(Gr_n(\mathbb{R});\mathbb{Z}/2\mathbb{Z})$. This is by the canonical CW decomposition of $Gr_n(\mathbb{R})$. Since the number of $m$-cells are equal to the number of sequences $1\le\sigma_1\le\dots\le\sigma_n\le m$. But the number of such sequences are exactly equal to the number of monomials formed by $w_i(\eta^n)$ of degree $m$. Thus $w_i(\eta^n)$ freely generate $H^*(Gr_n(\mathbb{R});\mathbb{Z}/2\mathbb{Z})$. \\
		Finally, we are ready to prove the uniqueness. Firstly, since the canonical bundle $\gamma_1^1$ of $\mathbb{RP}^1$ has Stiefel-Whitney class $1+a$, the canonical bundle $\gamma^1$ of $\mathbb{RP}^\infty$ has Stiefel-Whitney class $1+a$ by the canonical inclusion. Then, let the $n$-fold product $\xi=\gamma^1\times\cdots\times\gamma^1$, we have $w(\xi)=(1+a_1)\cdots(1+a_n)$. Now, since we have a bundle map $\xi\to\eta^n$, since $H^*(Gr_n(\mathbb{R});\mathbb{Z}/2\mathbb{Z})$ maps injectively into $H^*(\mathbb{RP}^\infty\times\cdots\times\mathbb{RP}^\infty;\mathbb{Z}/2\mathbb{Z})$, we have $w(\eta^n)$ unique. Thus by \cref{Real_Grassmannian}, we obtain the results. 
	\end{proof}

	Moreover, with the help of the Grassmannians, we can prove the following tensor formula. 
	\begin{thm}\label{SW_tensor}
		For $\xi$ an $m$-dimensional bundle and $\eta$ an $n$-dimensional bundle over $B$, the Stiefel-Whitney class $w(\xi\otimes\eta)$ is given by a universal formula
		\[ w(\xi\otimes\eta)=p_{m,n}(w_1(\xi),\dots,w_m(\xi),w_1(\eta),\dots,w_n(\eta)) \]
		Where $p_{m,n}$ can be characterized by the following identity: 
		\[ p_{m,n}(\sigma_1,\dots,\sigma_m,\sigma'_1,\dots,\sigma'_n)=\prod_{i=1}^m\prod_{j=1}^n(1+t_i+t'_j) \]
		where $\sigma$ and $\sigma'$ are the elementary symmetric polynomials of $t_i$ and $t'_j$. 
	\end{thm}
	\begin{proof}
		It suffices to check the formula on the universal case: notice the operation of forming the tensor product bundle is characterized by a map $Gr_m(\mathbb{R})\times Gr_n(\mathbb{R})\to Gr_{mn}(\mathbb{R})$. However, recall that we have a natural map inducing an injection on mod-2 cohomology $Gr_1(\mathbb{R})^k\to Gr_k(\mathbb{R})$, given by the classifying map of the $k$-fold direct sum of the canonical bundle. Therefore, we may consider the following diagram 
		% https://q.uiver.app/#q=WzAsNCxbMCwwLCJHcl8xKFxcbWF0aGJie1J9KV57bStufSJdLFsxLDAsIkdyX20oXFxtYXRoYmJ7Un0pXFx0aW1lcyBHcl9uKFxcbWF0aGJie1J9KSJdLFswLDEsIkdyXzEoXFxtYXRoYmJ7Un0pXnttbn0iXSxbMSwxLCJHcl97bW59KFxcbWF0aGJie1J9KSJdLFswLDFdLFsxLDNdLFswLDJdLFsyLDNdXQ==
		\[\begin{tikzcd}
			{Gr_1(\mathbb{R})^{m+n}} & {Gr_m(\mathbb{R})\times Gr_n(\mathbb{R})} \\
			{Gr_1(\mathbb{R})^{mn}} & {Gr_{mn}(\mathbb{R})}
			\arrow["\varphi_{m+n}", from=1-1, to=1-2]
			\arrow["\otimes"', from=1-1, to=2-1]
			\arrow["\otimes", from=1-2, to=2-2]
			\arrow["\varphi_{mn}"', from=2-1, to=2-2]
		\end{tikzcd}\]
		where the left map is again the tensor map (since the tensor of vector bundles commutes with direct sums, the tensor product of direct sums of line bundles decomposes into direct sums of line bundles). \\
		Thus it suffices to determine the map on the left. \\
		We first consider the case where $m=n=1$. Recall that $H^*(Gr_1(\mathbb{R})^2;\mathbb{Z}/2\mathbb{Z})=\mathbb{Z}/2\mathbb{Z}[t_1, t_2]$. Therefore $\otimes^*:H^*(Gr_1(\mathbb{R}))\to H^*(Gr_1(\mathbb{R})^2)$ is determined by $\otimes^*(t)=at_1+bt_2$ for some constant $a$ and $b$. In this case, since the tensor product of a bundle with the 1-dimensional trivial bundle is itself, we have $a=b=1$. Explicitly, we have proven that for two line bundles $\xi$ and $\eta$, $w_1(\xi\otimes\eta)=w_1(\xi)+w_1(\eta)$. \\
		Then we consider the map $\otimes^*:H^*(Gr_1(\mathbb{R})^{mn})\to H^*(Gr_1(\mathbb{R})^{m+n})$. The left is a polynomial ring with generators $t_{i,j}$, and the right is a polynomial ring with generators $s_i$ and $s'_j$. This map is given by $t_{i,j}\mapsto s_i+s'_j$, by the case of line bundles and the product formula of direct sum bundles. \\
		Finally for the general case, recall that under the inclusion $H^*(Gr_n(\mathbb{R}))\to H^*(Gr_1(\mathbb{R})^n)$, the image of $w_k(\eta^n)$ is $\sigma_k$. Therefore, we have
		\begin{align*}
			&\otimes^*(w(\eta^{mn})) \\
			=&\otimes^*(1+w_1(\eta^{mn})+\cdots+w_{mn}(\eta^{mn})) \\
			=&\otimes^*(1+\sigma_1+\cdots+\sigma_{mn}) \\
			=&\otimes^*(\prod(1+t_{ij})) \\
			=&\prod(1+t_i+t'_j)
		\end{align*}
		This implies the general case. 
	\end{proof}
	\begin{remark}
		The method used by passing from the case of $Gr_n$ to $Gr_1$ is called the splitting principle. We will repeatedly use this method throughout this part of the note. 
	\end{remark}

\section{Euler classes}

	Before we finally go into the non-trivial results of Stiefel-Whitney classes, we give a closer inspect to $\mathbb{Z}$-oriented vector bundles. 
	\begin{defn}
		For an $\mathbb{Z}$-oriented $n$-plane bundle $\xi$, by \cref{fund_class}, the fundamental class $u\in H^n(E,E_0;\mathbb{Z})$ always exists. Then we define the Euler class of $\xi$, $e(\xi)$, to be image $(\pi^*)^{-1}i^*(u)\in H^n(B;\mathbb{Z})$, where $i^*$ is the canonical map $H^n(E,E_0;\mathbb{Z})\to H^n(E,\mathbb{Z})$, and $\pi^*$ is the projection isomorphism $H^n(B;\mathbb{Z})\to H^n(E;\mathbb{Z})$. 
	\end{defn}

	Unlike the $\mathbb{Z}/2\mathbb{Z}$ case, the sign of orientation matters. 
	\begin{coro}
		For an orientation preserving bundle map $f:\xi\to\eta$ ($\xi$ and $\eta$ are always assumed to have the same dimension), $f^*w(\eta)=w(\xi)$. \\
		If the choice of orientation of $\xi$ is inversed, then the Euler class changes sign. 
	\end{coro}
	\begin{proof}
		Direct by definition. 
	\end{proof}

	\begin{coro}
		For any odd dimensional bundle $\xi$, we have $2e(\xi)=0$. 
	\end{coro}
	\begin{proof}
		By the Thom isomorphism (\cref{Thom_isom}) $\varphi(x)=p^*(x)\smile u$ maps $e(\xi)$ to the class $u\smile u$. Thus $u$ having odd dimension implies $2u\smile u=0$, thus $2e(\xi)=0$. 
	\end{proof}

	Naturally, the Euler class relates to the top Stiefel-Whitney class. 
	\begin{prop}
		The natural homomorphism $H^n(B;\mathbb{Z})\to H^n(B;\mathbb{Z}/2\mathbb{Z})$ maps the Euler class (if it exists) to the top Stiefel-Whitney class. 
	\end{prop}
	\begin{proof}
		By the property of \cref{Steenrod_op}, we have $Sq^n(u)=u\smile u$ since $u$ has dimension $n$. Now $e(\xi)=\varphi^{-1}(u\smile u)$, hence $e(\xi)$ maps to the class $\varphi^{-1}(u\smile u)$ in mod $2$ cohomology, which is exactly $w_n(\xi)$. 
	\end{proof}

	\begin{prop}\label{Euler_prod}
		The Euler class of a product is given by $e(\xi\times\xi')=e(\xi)\times e(\xi')$. Similarly the Euler class of a Whitney sum is given by $e(\xi\oplus\xi')=e(\xi)\smile e(\xi')$. 
	\end{prop}
	\begin{proof}
		Suppose $\xi$ has dimension $m$ and $\xi'$ has dimension $n$. We first note that for the fundamental class we have $u(\xi\times\xi')=(-1)^{mn}u(\xi)\times u(\xi')$. Then considering the restriction homomorphism $H^{m+n}(E\times E',(E\times E')_0)\to H^{m+n}(E\times E')\simeq H^{m+n}(B\times B')$ gives $e(\xi\times\xi')=(-1)^{mn}e(\xi)\times e(\xi')$. \\
		Now, the Whitney sum case follows from pulling the product bundle along the diagonal $B\to B\times B$, hence $e(\xi\oplus\xi')=e(\xi)\smile e(\xi')$ by the definition of cross products. 
	\end{proof}

	\begin{coro}
		If a bundle has a nowhere vanishing cross section, then the Euler class is $0$. 
	\end{coro}
	\begin{proof}
		$e(\xi)=e(\varepsilon^1)\smile e(\varepsilon^{1\bot})$. 
	\end{proof}


\chapter{Computations in Manifolds}

\section{Some common bundles of a manifold}

	We begin with the most canonical bundle associated to any manifold: the tangent bundle. But before we really start anything, we need some lemmas on normal bundles and tangent bundles. 
	\begin{lemma}
		The orientation of the tangent bundle $\tau_M$ is equivalent to an orientation on $M$. 
	\end{lemma}
	\begin{proof}
		Recall the orientation of a manifold is a choice of a preferred generator $\mu_x\in H_n(M,M-x;A)$ compatible in the sense that $\mu_x$ and $\mu_y$ in some neighborhood $N$ homeomorphic to $\mathbb{R}^n$ comes from the same class in $H_n(M,M-N;A)$. \\
		Similarly, the orientation of the bundle can be given by choosing a consistent generator $\mu_x\in H_n(\tang_xM,\tang_xM-0;A)$ (under the natural pairing of the top cohomology and homology classes of spheres). But by the preceding corollary, we have $H_n(M,M-x;A)\simeq H_n(\tang_xM,\tang_xM-0;A)$, which gives the results. 
	\end{proof}
	\begin{lemma}
		The normal bundle $\nu^n$ associated with the diagonal embedding of $M$ into $M\times M$ is canonically isomorphic to the tangent bundle of $M$. 
	\end{lemma}
	\begin{proof}
		We consider the tangent space $(u,v)\in \tang_xM\times \tang_xM\simeq$ is tangent to $\Delta(M)$ iff $u=v$ and normal to $\Delta(M)$ iff $u=-v$. Thus $(x,v)\mapsto((x,x),(-v,v))$ gives an isomorphism of bundles. 
	\end{proof}

	The next result arises from differential topology, and is crucial to some of the later results. 
	\begin{thm}\label{Tube_Nbd}
		For an embedding of manifolds $f:M\to N$, there is an neighborhood of $M$, say $M'\subset N$, such that $M'$ is diffeomorphic to the total space of the normal bundle of $f(M)$ such that $x\in M$ has the image $f(x)$ the zero vector in the bundle under the identification $M'\simeq\nu_M$. 
	\end{thm}

	\begin{coro}
		If $f(M)$ is closed in $N$, then the cohomology ring $H^*(E,E_0;A)$ associated to the normal bundle of $M$ in $N$ is canonically isomorphic to $H^*(N,N-M;A)$. The same property applies to the homology. 
	\end{coro}
	\begin{proof}
		By excision, we have $H^*(N,N-M;A)\simeq H^*(M',M'-M;A)\simeq H^*(E,E_0;A)$. 
	\end{proof}

	With the help of the tubular neighborhood lemma, we have a characterization of the Euler class of the normal bundle, and, consequently, the top Stiefel-Whitney class. 
	\begin{prop}\label{Fund_Euler}
		If $M$ is embedded as a closed submanifold of $N$ and the normal bundle $\nu^k$ is oriented0, the composition of the restriction homomorphisms 
		\[ H^k(E,E_0;\mathbb{Z})\to H^k(N,N-M;\mathbb{Z})\to H^k(N;\mathbb{Z})\to H^k(M;\mathbb{Z}) \]
		sends the fundamental class $u$ to the Euler class $e(\nu^k)$. Consequently, the modulo $2$ fundamental class is send to the top Stiefel-Whitney class $w_k(\nu^k)$. 
	\end{prop}
	\begin{proof}
		This is by the definition of the Euler class, and the top Stiefel-Whitney class case is completely similar. 
	\end{proof}

	\begin{coro}
		If an $n$-dimensional manifold $M$ is embedded as a closed submanifold of $\mathbb{R}^{n+k}$, we have the top Stiefel-Whitney class $w_k(\nu^k)=0$. In the oriented case $e(\nu^k)=0$. 
	\end{coro}
	\begin{proof}
		$H^k(N;\Lambda)=0$. 
	\end{proof}

\section{The diagonal cohomology class}

	\begin{defn}
		For a fixed coefficient ring $\Lambda$, we assume that an $n$-dimensional manifold $M$ is $\Lambda$ oriented. We define the cohomology class $u'\in H^n(M\times M,M\times M-\Delta(M);\Lambda)$ to be the fundamental class of the normal bundle of the diagonal embedding $\Delta:M\to M\times M$ (via the identification $H^n(M\times M,M\times M-\Delta(M);\Lambda)\simeq H^n(E,E_0;\Lambda)$ where $E$ is the total space of the normal bundle). \\
		Note that here this is well-defined since when $M$ is $\Lambda$-oriented, the tangent bundle of $M$ is $\Lambda$-oriented, and under the isomorphism $\tau_M\simeq\nu^n$ in the diagonal embedding, the fundamental class of the normal bundle is well-defined. 
	\end{defn}

	We have a more explicit characterization of the fundamental class. 
	\begin{prop}\label{Fund_char}
		For an oriented $n$-dimensional manifold $M$, each $x\in M$, the cohomology class $H^n(M,M-x;\Lambda)$ has a preferred generator $u_x$ given by the perfect pairing $\langle u_x,\mu_x\rangle=1$. Then, the diagonal cohomology class can be characterized as follows: under the inclusion $j_x:(M,M-x)\inj(M\times M,M\times M-\Delta(M))$ given by $y\mapsto(x,y)$, we have $j_x^*(u')=u_x$ for all $x$. 
	\end{prop}
	\begin{proof}
		By its definition, it admits the following unique characterization: for any $x$ and any small neighborhood $N$ of $0$ in the tangent space $\tang_xM$, consider the embedding $(N,N-0)\to(M\times M,M\times M-\Delta(M))$, defined by the tubular neighborhood lemma. Then the induced cohomology morphism must map $d_M$ to the generator $u_x$. The rest follows from the embedding $(N,N-0)\to(M,M-X)$. 
	\end{proof}
	
	\begin{defn}
		From the fundamental class, we define the diagonal cohomology class as the image of the restriction homomorphism $H^n(M\times M,M\times M-\Delta(M))\to H^n(M\times M)$ of the fundamental class $u'$. We denote this class by $d_M$. 
	\end{defn}
	
	As its name, it behaves somehow like a diagonal. 
	\begin{lemma}
		For any class $a\in H^*(M)$, we have $(a\times 1)\smile d_M=(1\times a)\smile d_M$. 
	\end{lemma}
	\begin{proof}
		Let $N$ be a tubular neighborhood of the diagonal $\Delta(M)$ in $M\times M$, hence $\Delta(M)$ is a deformation retract of $N$. Now the projection $p_1,p_2:M\times M\to M$ coincides on $\Delta(M)$, hence the restriction $p_1|_N$ and $p_2|_N$ are homotopic to each other. Thus we have $p_1^*(a)=a\times 1$ and $p_2^*(a)=1\times a$ has the same image under the restriction homomorphism $\iota^*:H^i(M\times M)\to H^i(N)$. This gives rise to the following commutative diagram: 
		% https://q.uiver.app/#q=WzAsNCxbMCwwLCJIXmkoTVxcdGltZXMgTSkiXSxbMSwwLCJIXmkoTikiXSxbMCwxLCJIXntpK259KE1cXHRpbWVzIE0sTVxcdGltZXMgTS1cXERlbHRhKE0pKSJdLFsxLDEsIkhee2krbn0oTixOLVxcRGVsdGEoTSkpIl0sWzAsMV0sWzEsMywiXFxzbWlsZSBcXGlvdGFeKnUnIl0sWzAsMiwiXFxzbWlsZSB1JyIsMl0sWzIsMywiXFxzaW1lcSIsMl1d
		\[\begin{tikzcd}
			{H^i(M\times M)} & {H^i(N)} \\
			{H^{i+n}(M\times M,M\times M-\Delta(M))} & {H^{i+n}(N,N-\Delta(M))}
			\arrow[from=1-1, to=1-2]
			\arrow["{\smile \iota^*u'}", from=1-2, to=2-2]
			\arrow["{\smile u'}"', from=1-1, to=2-1]
			\arrow["\simeq"', from=2-1, to=2-2]
		\end{tikzcd}\]
		Thus we have $1\times a$ and $a\times 1$ are equal when we pass to the bottom right cohomology group, hence by the bottom equivalence, we have $(1\times a)\smile u'=(a\times 1)\smile u'$, and the restriction to $H^{i+n}(M\times M)$ gives $(1\times a)\smile d_M=(a\times 1)\smile d_M$. 
	\end{proof}
	
	\begin{defn}
		We recall that we have a slant product operation $E^{p+q}(X\times Y;\Lambda)\otimes E_q(Y;\Lambda)$ \\
		$\to E^p(X;\Lambda)$ by the following construction. Recall that all Eilenberg-Maclane spectra are ring spectrums. Then the information on the left is a morphism $f:X\wedge Y\to E$ and a morphism $g:\mathbb{S}\to E\wedge Y$. We may construct a map $X\to E$ as follows: 
		\[ X\xrightarrow{\simeq}X\wedge\mathbb{S}\xrightarrow{1\wedge g}X\wedge E\wedge Y\xrightarrow{1\wedge c}X\wedge Y\wedge E\xrightarrow{f\wedge 1}E\wedge E\xrightarrow{\mu}E \]
		where $\mu$ denotes the multiplication of the ring spectrum (for the explicit definition of smash products see part 2). Here we may substitute $E$ for any Eilenberg-Maclane spectra to get the desired definition. 
	\end{defn}
	\begin{lemma}
		For $a\in E^p(X)$, $b\in E^q(Y)$ and $c\in E_q(Y)$, we have $(a\times b)/c=a\langle b,c\rangle$. 
	\end{lemma}
	\begin{proof}
		Now we have $a\in[X,E]_{-p}$, $b\in[Y,E]_{-q}$, then $a\times b$ is by definition induced by the following map: 
		\[ X\wedge Y\xrightarrow{a\wedge b}E\wedge E\xrightarrow{\mu}E \]
		Then $(a\times b)/c$ is induced by the composition 
		\[ X\xrightarrow{\simeq}X\wedge\mathbb{S}\xrightarrow{1\wedge c}X\wedge E\wedge Y\xrightarrow{1\wedge c}X\wedge Y\wedge E\xrightarrow{a\times b\wedge 1}E\wedge E\xrightarrow{\mu}E \]
		Then we notice that the morphism the $X$ part is always the identity until the final part. This implies the result by noticing the pairing $\langle-,-\rangle$ is a special case of the slant product. 
	\end{proof}

	By virtue of this definition and the above lemma, we have the following property of the diagonal class: 
	\begin{prop}
		Suppose $M$ is compact and $\Lambda$-oriented so that we have the fundamental homology class $\mu$, then we have $d_M/\mu=1\in H^0(M;\Lambda)$. 
	\end{prop}
	\begin{proof}
		We consider the image of $d_M/\mu$ under the restriction morphism $H^0(M)\to H^0(x)$ for any $x$. We consider the commutative diagram % https://q.uiver.app/#q=WzAsNCxbMCwwLCJIXm4oTVxcdGltZXMgTSkiXSxbMSwwLCJIXjAoTSkiXSxbMCwxLCJIXm4oeFxcdGltZXMgTSkiXSxbMSwxLCJIXjAoeCkiXSxbMCwxLCIvXFxtdSJdLFswLDJdLFsyLDMsIi9cXG11IiwyXSxbMSwzXV0=
		\[\begin{tikzcd}
			{H^n(M\times M)} & {H^0(M)} \\
			{H^n(x\times M)} & {H^0(x)}
			\arrow["{/\mu}", from=1-1, to=1-2]
			\arrow[from=1-1, to=2-1]
			\arrow["{/\mu}"', from=2-1, to=2-2]
			\arrow[from=1-2, to=2-2]
		\end{tikzcd}\]
		Note that the left arrow sends $d_M$ to $1\times(i_x,1_M)^*(d_M)$, where $(i_x,1_M):M\to M\times M$ by $y\mapsto(x,y)$. Thus $i_x^*(d_M/\mu)=1\times(i_x,1_M)^*(d_M)/\mu=1\times\langle(i_x,1_M)^*(d_M),\mu\rangle\in H^0(x)$. \\
		Then, recall that the fundamental homology class $\mu$ is characterized by the fact that the image of $\mu$ under the homomorphism $H^n(M)\to H^n(M,M-x)$ is the preferred generator $\mu_x$. We now prove $\langle(i_x,1_M)^*(d_M),\mu\rangle=1$. We recall that $d_M$ comes from the fundamental class $u'\in H^n(M\times M,M\times M-\Delta(M))$. Now, by the following commutative square 
		% https://q.uiver.app/#q=WzAsNCxbMCwwLCJIXm4oTVxcdGltZXMgTSxNXFx0aW1lcyBNLVxcRGVsdGEoTSkpIl0sWzEsMCwiSF5uKE1cXHRpbWVzIE0pIl0sWzAsMSwiSF5uKE0sTS14KSJdLFsxLDEsIkhebihNKSJdLFswLDFdLFswLDIsImpfeCIsMl0sWzIsM10sWzEsMywiKGlfeCwxX00pXioiXV0=
		\[\begin{tikzcd}
			{H^n(M\times M,M\times M-\Delta(M))} & {H^n(M\times M)} \\
			{H^n(M,M-x)} & {H^n(M)}
			\arrow[from=1-1, to=1-2]
			\arrow["{j_x}"', from=1-1, to=2-1]
			\arrow[from=2-1, to=2-2]
			\arrow["{(i_x,1_M)^*}", from=1-2, to=2-2]
		\end{tikzcd}\]
		we have $(i_x,1_M)^*(d_M)=j_x^*(u')|_M$. But by \cref{Fund_char}, we have $\langle j_x^*(u'),\mu_x\rangle=1$, which is equal to $\langle j_x^*(u')|_M,\mu\rangle$. Thus we have $(d_M/\mu)|_x=1$ for all $x$, thus $d_M/\mu=1$. 
	\end{proof}

	\begin{thm}\label{Duality}
		Suppose $M$ is $\Lambda$-oriented, where $H^*(M)$ is free with basis $b_1,\dots,b_r\in H^*(M)$, then there exists a dual basis $b_1^\#,\dots,b_r^\#\in H^*(M)$ such that $\langle b_i\smile b_j^\#,\mu\rangle=\delta_{ij}$. \\
		Moreover, we have the following formula for the diagonal class
		\[ d_M=\sum_{i=1}^r(-1)^{\dim b_i}b_i\times b_i^\# \]
	\end{thm}
	\begin{proof}
		By the Kunneth formula, we have $d_M=b_1\times c_1+\cdots+b_r\times c_r$, where $\dim b_i+\dim c_i=n$. We apply the operation $/\mu$ to $(a\times 1)\smile d_M=(1\times a)\smile d_M$, and then we have 
		$((a\times 1)\smile d_M)/\mu=a\smile(d_M/\mu)=a$. On the right side, we substitute $\sum b_i\times c_i$ for $d_M$, then we have 
		\[ ((1\times a)\smile d_M)/\mu=\sum(-1)^{\dim a\dim b_i}(b_i\times(a\smile c_i))/\mu=\sum(-1)^{\dim a\dim b_i}b_i\langle a\smile c_i\rangle \]
		Thus the last expression is equal to $a$. Then since $a$ is an arbitrary cohomology class, we let it be $b_j$ and obtain that the coefficient $(-1)^{\dim a\dim b_i}\langle b_j\smile c_i,\mu\rangle$ must be $\delta_{ij}$. Thus we may let $b_i^\#=(-1)^{\dim b_i}c_i$, then we obtain the conclusion.  
	\end{proof}

	\begin{coro}\label{Euler_class_char}
		If $M$ is compact oriented, then we have $\chi(M)=\langle e(\tau_M),\mu\rangle$. 
	\end{coro}
	\begin{proof}
		We have $e(\tau_M)=\Delta^*(d_M)$ by \cref{Fund_Euler}. We may pass to rational coefficients and apply \cref{Duality}, then we have
		\[ d_M=\sum(-1)^{\dim b_i}b_i\times b_i^\# \]
		thus we have
		\[ e(\tau_M)=\sum(-1)^{\dim b_i}b_i\smile b_i^\# \] 
		By the characterization of the dual basis, we have the formula 
		\[ \langle e(\tau_M),\mu\rangle=\sum(-1)^{\dim b_i}=\chi(M) \]
	\end{proof}

\section{Wu class}

	\begin{defn}\label{Wu_class}
		We have an additive homomorphism $x\mapsto\langle Sq^k(x),\mu\rangle$ from \\
		$H^{n-k}(X;\mathbb{Z}/2\mathbb{Z})\to\mathbb{Z}/2\mathbb{Z}$. By \cref{Duality}, we have a cohomology class $v_k\in H^k(X;\mathbb{Z}/2\mathbb{Z})$ such that we have $\langle v_k\smile x,\mu\rangle=\langle Sq^k(x),\mu\rangle$ by letting $v_k$ be the sum of the dual basis of degree $k$. Actually, we may notice that we have $v_k\smile x=Sq^k(x)$. \\
		We define the \emph{total Wu class} $v\in H^\Pi(X)$ to be the formal sum $v=1+\cdots+v_n$. Clearly we have $\langle v\smile x,\mu\rangle=\langle Sq(x),\mu\rangle$. 
	\end{defn}

	\begin{lemma}
		We have $w_i(\tau_M)=Sq^i(d_M)/\mu$. 
	\end{lemma}
	\begin{proof}
		Firstly by \cref{Cons_SW}, we have $w_i(\tau_M)=\varphi^{-1}Sq^i(u)$, where $\varphi$ denotes the Thom isomorphism (in \cref{Thom_isom}), hence we have $Sq^i(u)=\varphi(w_i(\tau_M))=(\pi^*w_i(\tau_M))\smile u$. With the identification $H^*(E,E_0)\simeq H^*(M\times M,M\times M-\Delta(M))$, we have $Sq^i(u')=(w_i(\tau_M)\times 1)\smile u'$. Thus after passing to $H^*(M\times M)$, we have $Sq^i(d_M)=(w_i(\tau_M)\times 1)\smile d_M$. \\
		Therefore, by applying the slant product $/\mu$ on both sides of the equation, we have $Sq^i(d_M)/\mu=((w_i(\tau_M)\times 1)\smile d_M)/\mu=w_i\smile(d_M/\mu)=w_i$. 
	\end{proof}
	\begin{thm}
		For any manifold, the total Stiefel-Whitney class $w$ has a characterization $w(\tau_M)=Sq(v)$. 
	\end{thm}
	\begin{proof}
		We choose a basis $b_i$ for the modulo $2$ cohomology class of $M$ and a dual basis $b_i^\#$ by \cref{Duality}. Then we have $x=\sum b_i\langle x\smile b_i^\#,\mu\rangle$. Therefore by substituting $u$ for $x$ and applying the total Steenrod operation on both sides, we have
		\[ Sq(v)=\sum Sq(b_i)\langle Sq(b_i^\#),\mu\rangle=\sum(Sq(b_i)\times Sq(b_i^\#))/\mu=Sq(d_M)/\mu=w(M) \]
		by the preceding lemma. 
	\end{proof}


\chapter{Chern Classes and Pontrajagin Classes}

\section{A quick review}

	\begin{defn}
		An $n$-dimensional complex manifold $M$ is a second countable Hausdorff space where each point $x\in M$ such that it has a neighborhood homeomorphic to an open subset in $\mathbb{C}^n$, and the transitions between these open sets are given by holomorphic functions. 
	\end{defn}

	\begin{defn}
		A complex bundle over a space $B$ is a projection $p:E\to B$, such that for each $b\in B$, the fibre $\pi^{-1}(b)$ can be identified with a complex vector space $V$. Moreover, it need to satiesfy the local triviality condition: for any $b\in B$, there exists a neighborhood $U$, such that the fibre $\pi^{-1}(U)$ is homeomorphic to $U\times V$. Moreover, the inclusion of each fibre $\pi^{-1}(b)\inj\pi^{-1}(U)$ must be complex linear. 
	\end{defn}

	\begin{defn}
		A complex structure on a real bundle $\xi$ is a continuous bundle map $J:E(\xi)\to E(\xi)$ such that $J$ is $\mathbb{R}$-linear on each fibre and satisfies the identity $J^2(v)=-v$ for every vector $v$ in the fibre. 
	\end{defn}
	\begin{coro}
		If a $\mathbb{R}$-vector bundle admits a complex structure, then it has dimensions $2n$ for some $n$. Moreover, it can be regarded as a $n$-dimensional $\mathbb{C}$-vector bundle. 
	\end{coro}

	\begin{defn}
		A almost complex structure on a real manifold $M$ is a complex structure on its tangent bundle $\tau_M$. 
	\end{defn}

\section{Chern classes}

	\begin{lemma}
		For any complex vector bundle $\omega$, its underlying real vector bundle $\omega_\mathbb{R}$ has a preferred orientation. 
	\end{lemma}
	\begin{proof}
		For any complex vector space $V$, we choose its basis $v_1,\dots,v_n$. Then $v_1,iv_1,\dots,v_n,iv_n$ determins a real basis of the underlying real vector space $V_\mathbb{R}$. Also, the real determinant of this set of basis is always positive. Hence a complex vector space has a preferred orientation. We may then pass easily to complex vector bundles. 
	\end{proof}

	\begin{coro}
		For any complex vector bundle $\omega$ on $B$, its Euler class $e(\omega_\mathbb{R})\in H^{2n}(B;\mathbb{Z})$ is well-defined. 
	\end{coro}

	\begin{thm}\label{Gysin_seq}
		For a $n$-dimensional $\Lambda$-oriented bundle $\xi$, we have the following sequence: 
		\[ \cdots\lto H^{i-n}(B;\Lambda)\lto H^i(B;\Lambda)\lto H^i(E_0;\Lambda)\lto H^{i-n+1}(B;\Lambda)\lto\cdots \]
	\end{thm}
	\begin{proof}
		We already have the long exact sequence of cohomology
		\[ \cdots\lto H^i(E,E_0;\Lambda)\lto H^i(E;\Lambda)\lto H^i(E_0;\Lambda)\lto H^{i+1}(E,E_0;\Lambda) \]
		The rest follows from \cref{Thom_isom} by orientedness of $\xi$. 
	\end{proof}
	We now construct the Chern class of a complex bundle. 

	\begin{const}\label{Chern_const}
		For an $n$-dimensional complex bundle $\omega$ on $B$, we define its top Chern class $c_n(\omega)$ to be the Euler class $e(\omega_\mathbb{R})$. Then, we inductively define the lower classes. \\
		Note that firstly, for a $n$-dimensional bundle $\omega$, we have an $n-1$-dimensional bundle $\omega_0$ on $E_0(\omega)$. Namely, we associate each point $v_x\in E(\omega)$ to its orthogonal complement in $F_x$. Then, we define $c_i(\omega)=(\pi_0^*)^{-1}(c_i(\omega_0))$. Note that $\pi_0^*:H^{2i}(B)\to H^{2i}(E_0)$ is an isomorphism for $i<n$ by \cref{Gysin_seq}, so that the inverse image makes sense. \\
		We then define the formal sum $c(\omega)=1+c_1(\omega)+\cdots+c_n(\omega)$ to be the total Chern class of $\omega$. Note that this is a unit. 
	\end{const}

	\begin{lemma}\label{Chern_nat}
		For a bundle map $f_*:\omega\to\omega'$ over $f:B\to B'$, we have $c(\omega)=f^*(c(\omega'))$. 
	\end{lemma}
	\begin{proof}
		Directly by the naturality of the Euler class and the naturality of the construction $\omega_0$. 
	\end{proof}

	Similar to the real Grassmannian classifying real vector bundles, we have the complex Grassmannian classifies complex vector bundles. 
	\begin{thm}\label{Cplx_Grassmannian}
		For a paracompact topological space $X$, an $n$-dimensional $\mathbb{C}$-vector bundle $p:E\to X$ corresponds to a homotopy class of maps $f:X\to Gr_n(\mathbb{C})$. Moreover, this bundle is the pullback bundle of the canonical $n$-bundle $\eta^n$ over $Gr_n(\mathbb{C})$. 
	\end{thm}

	We now give another characterization of Chern classes. 
	\begin{thm}\label{Chern_gen}
		The cohomology ring $H^*(Gr_n(\mathbb{C});\mathbb{Z})$ is the polynomial ring over $\mathbb{Z}$ generated by the Chern classes $c_1(\eta^n),\dots,c_n(\eta^n)$. 
	\end{thm}
	\begin{remark}
		In the following parts, we have an analogous statement claiming that for suitable cohomology theory $E$, we have $E^*(Gr_n)\simeq(\pi_*(E))[[c_i]]$. Notice it is only in this part we adopt the convention that $H^*\simeq\bigoplus H^n$. In later parts, we will define $H^*\simeq\prod H^n$. 
	\end{remark}

	Before proving the theorem, we need some lemmas. 
	\begin{lemma}
		The cohomology ring $H^*(\mathbb{CP}^k;\mathbb{Z})$ is a truncated polynomial ring terminating in dimension $2k$, and generated by the first Chern class of the canoncial line bundle $c_1(\gamma_k^1)$. 
	\end{lemma}
	\begin{proof}
		We use the Gysin sequence (\cref{Gysin_seq}) on the canonical line bundle $\gamma_k^1$ of $\mathbb{CP}^k$. Using the definition of Chern classes, $c_1(\gamma_k^1)=e((\gamma_k^1)_\mathbb{R})$ since it is complex $1$-dimensional, we have 
		\[ \cdots\lto H^{i+1}(E_0;\mathbb{Z})\lto H^i(\mathbb{CP}^k;\mathbb{Z})\lto H^{i+2}(\mathbb{CP}^k;\mathbb{Z})\lto H^{i+2}(E_0;\mathbb{Z})\lto\cdots \]
		Now that the space $E_0$ is the set of all pairs $(\ell,v)$ such that $\ell$ is a $1$-dimensional subspace in $\mathbb{C}^{k+1}$, and $v\in\ell$ is non-zero. This can be identified with $\mathbb{C}^{k+1}-\{0\}$, hence is homotopic to $S^{2k+1}$. Thus for $0\le i\le 2k-2$, we have $H^{i+2}(E_0)=0$. Thus we have an isomorphism $\smile c_1:H^i(\mathbb{CP}^k;\mathbb{Z})\to H^{i+2}(\mathbb{CP}^k;\mathbb{Z})$. Beginning with $H^{-1}$ gives the result that all odd dimensional cohomology vanishes, and beginning with $H^0$ gives that all even dimensional cohomology is $\mathbb{Z}$. Finally, the isomorphism given by $\smile c_1$ gives that $c_1$ generates the cohomology freely. 
	\end{proof}
	\begin{coro}
		$H^*(Gr_1(\mathbb{C});\mathbb{Z})$ is the free polynomial over $\mathbb{Z}$ generated by $c_1(\eta^1)$. 
	\end{coro}
	\begin{proof}
		The above argument applies to the case $k=\infty$. 
	\end{proof}
	
	Now we are ready to prove the theorem. 
	\begin{proof}
		We prove by induction on $n$. Note that the case $n=1$ is already solved in the preceding corollary. \\
		Then, consider the map $f:E_0(\eta^n)\to Gr_{n-1}(\mathbb{C})$ defined by the following process: for a point $(V,v)\in E_0$, where $V$ is an $n$-subspace in $\mathbb{C}^\infty$ and $v$ a point in it, we define $f(V,v)=V\cap v^\bot$, which is an $n-1$-subspace in $\mathbb{C}^\infty$. We note here that $f$ can also be defined equivalently as the map classifying the orthogonal $n-1$-bundle on $E_0$. \\
		We prove $f^*:H^*(Gr_{n-1}(\mathbb{C});\mathbb{Z})\to H^*(E_0;\mathbb{Z})$ is an isomorphism. We would need to take the orthogonal complement of a subspace, so we would first restrict $f$ to $f_N:E_0(\eta_N^n)\to Gr_{n-1}^N(\mathbb{C})$, where $Gr_m^n$ denotes the manifold of all $m$ subspace in $\mathbb{C}^n$, and $E_0(\eta_N^n)$ is the subbundle of $E_0(\eta^n)$ over $Gr_n^N(\mathbb{C})$. We note that we have a bundle $\omega_N^{N-n+1}$ over $Gr_{n-1}^N(\mathbb{C})$ where each fibre can be identified with the orthogonal complement of the $n-1$-subspace corresponding to a point in $Gr_{n-1}^N(\mathbb{C})$. This bundle can be identified with $f_N$ by definition. Using the Gysin sequence of this bundle, we have $f_N$ induces cohomology isomorphisms in dimension $\le 2(N-n)$. Using cellular filtration we have $f$ induces cohomology isomorphism in all dimensions. \\
		Then we may substitute $Gr_{n-1}(\mathbb{C})$ for $E_0$ in the Gysin sequence: 
		\[ \cdots\lto H^i(Gr_n)\lto H^{i+2n}(Gr_n)\lto H^{i+2n}(Gr_{n-1})\lto H^{i+1}(Gr_n)\lto\cdots \]
		We prove the homomorphism $\lambda:H^{i+2n}(Gr_n)\to H^{i+2n}(Gr_{n-1})$ sends the Chern class $c_i(\eta^n)$ to $c_i(\eta^{n-1})$. By the equivalent definition of $f:E_0\to Gr_{n-1}(\mathbb{C})$, we have $f^*(c_i(\eta^{n-1}))=c_i(\eta_0^n)$, and $\pi_0^*(c_i(\eta^n))=c_i(\eta_0^n)$ is by \cref{Chern_const}. \\
		It now remains to check the algebraic independence. We now apply the induction hypothesis. Since the map in the sequence $H^{i+2n}(Gr_n)\to H^{i+2n}(Gr_{n-1})$ is surjective, it splits into the following short exact sequences
		\[ 0\lto H^i(Gr_n)\xrightarrow{\smile c_n} H^{i+2n}(Gr_n)\xrightarrow{\lambda} H^{i+2n}(Gr_{n-1})\lto 0 \]
		By this sequence, we prove the algebraic independence. For $x\in H^{i+2n}(Gr_n)$, $\lambda(x)$ can be uniquely expressed as a polynomial, say $\lambda(x)=p(c_1(\eta^{n-1}),\dots,c_{n-1}(\eta^{n-1}))$, hence \\
		$x-p(c_1(\eta^n),\dots,c_{n-1}(\eta^n))$ is in the kernel of $\lambda$. Thus we have $x-p(c_1(\eta^n),\dots,c_{n-1}(\eta^n))=y\smile c_n(\eta^n)$ for some $y\in H^i(Gr_n)$. Now since the dimension of $y$ is lower then $x$, we apply another induction on the dimension and we may now suppose $y$ is uniquely a polynomial $q(c_1(\eta^n),\dots,c_n(\eta^n))$. Now that $x=p(c_1(\eta^n),\dots,c_{n-1}(\eta^n))+c_n(\eta^n)q(c_1(\eta^n),\dots,c_n(\eta^n))$, by applying $\lambda$ we see $p$ is unique, then $q$ is obviously unique. 
	\end{proof}

	Similar to the product formula of the Stiefel-Whitney class, we have the following formula
	\begin{thm}\label{Chern_prod}
		For two bundles $\omega$, $\varphi$ over $B$, we have $c(\omega\oplus\varphi)=c(\omega)c(\varphi)$. 
	\end{thm}

	\begin{lemma}
		There exists one and only polynomial $p_{m,n}=p_{m,n}(c_1,\dots,c_m;c'_1,\dots,c'_n)$ with the following identity $c(\omega\oplus\varphi)=p_{m,n}(c_1(\omega),\dots,c_m(\omega);c_1(\varphi),\dots,c_n(\varphi))$ for all $m$-bundle $\omega$ and $n$-bundle $\varphi$ over a common paracompact base space $B$. 
	\end{lemma}
	\begin{proof}
		Note the space $Gr_m\times Gr_n$ classifies a pair consisting of an $m$-bundle and an $n$-bundle, where the universal bundle is constructed by the Whitney sum of the pullbacks of the two universal bundles, denoted by $\eta_1^m\oplus\eta_2^n$, or the product bundle. \\
		Now, since $Gr_n$ has free abelian cohomology, we have the Kunneth formula available $H^*(Gr_m)\otimes H^*(Gr_n)\simeq H^*(Gr_m\times Gr_n)$. Therefore $H^*(Gr_m\times Gr_n)$ is a polynomial ring over $\mathbb{Z}$ on the generators $c_i(\eta_1^m)$ and $c_j(\eta_2^n)$ for $1\le i\le m$ and $1\le j\le n$. Now for $\omega$ and $\varphi$ on $B$, it is classified by a map $f:B\to Gr_m\times Gr_n$. Hence we may take $p_{m,n}$ to be the expression of $c(\eta_1^m\oplus\eta_2^n)=p_{m,n}(c_1(\omega),\dots,c_n(\omega);c_1(\varphi),\dots,c_n(\varphi))$. 
	\end{proof}

	\begin{lemma}
		$p_{m,n}$ is given by $(1+c_1+\cdots+c_m)(1+c'_1+\cdots+c'_n)$. 
	\end{lemma}
	\begin{proof}
		We proceed by induction on $m$ and $n$. Suppose inductively that $c(\eta_1^{m-1}\oplus\eta_2^n)$ is equal to $(1+c_1+\cdots+c_{m-1})(1+c'_1+\cdots+c'_n)$. Then for the bundle $\eta_1^{m-1}\oplus\varepsilon^1\oplus\eta_2^n$ over $Gr_{m-1}\times Gr_n$, we have 
		\[ c(\eta_1^{m-1}\oplus\varepsilon^1\oplus\eta_2^n)=c(\eta_1^{m-1}\oplus\eta_2^n)=p_{m,n}(c_1(\eta_1^{m-1}),\dots,c_m(\eta_1^{m-1})=0;c_1(\eta_2^n),\dots,c_n(\eta_2^n)) \]
		By induction hypothesis, we have 
		\[ p_{m,n}(c_1,\dots,c_m=0;c'_1,\dots,c'_n)=(1+c_1+\cdots+c_{m-1})(1+c'_1+\cdots+c'_n) \]
		By applying a similar induction on $n$, we have
		\[ p_{m,n}(c_1,\dots,c_m;c'_1,\dots,c'_n)=(1+\cdots+c_m)(1+\cdots+c'_n)+uc_mc_n \]
		By considering the highest dimension we have $u$ has dimension $0$, or the total Chern class will have a element of dimension greater than $2m+2n$. But $u=0$ since the top Chern class can be identified with the Euler class, which satisfies the product formula. 
	\end{proof}
	Now \cref{Chern_prod} is a direct corollary. \\
	With the help of this theorem, we may prove the following analogous statement in $Gr_n(\mathbb{C})$ as in $Gr_n(\mathbb{R})$. 
	\begin{thm}
		Consider the product of the pullback bundles on $Gr_1(\mathbb{C})^n$, $\xi=\eta^1_1\oplus\cdots\oplus\eta^1_n$. This bundle is classified by a map $\varphi:Gr_1(\mathbb{C})^n\to Gr_n(\mathbb{C})$. Then $\varphi^*(c_k(\eta^n))=\sigma_k$, the $k$-th elementary symmetric polynomial of $c_1(\eta^1_1),\dots,c_1(\eta^1_n)$. 
	\end{thm}
	\begin{proof}
		We have 
		\[ c(\xi)=c(\eta^1_1)\cdots c(\eta^1_n)=\sum\sigma_k(c_1(\eta^1_1),\dots,c_1(\eta^1_n)) \]
		The rest follows by \cref{Chern_nat}. 
	\end{proof}

	Similar to the definition of Euler classes, some issues with orientations arises in Chern classes. 
	\begin{defn}
		For a complex $n$-bundle $\omega$, we define its conjugate bundle $\bar{\omega}$ to be the bundle with the same underlying real bundle $\omega_\mathbb{R}$, and with the opposite complex structure, that is, $\bar{J}(x)=-J(x)$. Thus the identity map $f:E(\omega)\to E(\bar{\omega})$ is conjugate linear, explicitly $f(\lambda e)=\bar{\lambda}f(e)$. 
	\end{defn}

	\begin{prop}\label{Chern_ori}
		We have $c(\bar{\omega})=1-c_1(\omega)+\cdots+(-1)^nc_n(\omega)$. 
	\end{prop}
	\begin{proof}
		We prove by induction, hence it suffices to verify the top Chern class. But this follows from the vector space case by the following argument: \\
		If $\{v_j,iv_j\}_{1\le j\le n}$ is chosen as a basis for $F$, then $\{v_j,-iv_j\}_{1\le j\le n}$ is the preferred basis for $\bar{F}$. Then the orientation is reversed for $n$ times, thus for the top class we have $c_n(\omega)=(-1)^nc_n(\bar{\omega})$. 
	\end{proof}
	Finally, analogous to \cref{SW_tensor}, we have the following theorem. 
	\begin{thm}\label{Chern_tensor}
		For an $m$-dimensional complex vector bundle $\xi$ and an $n$-dimensional bundle $\eta$, we have $c(\xi\otimes\eta)=p_{m,n}(c_1(\xi),\dots,c_m(\xi),c_1(\eta),\dots,c_n(\eta))$, where $p_{m,n}$, as before, is given by
		\[ p_{m,n}(\sigma_1,\dots,\sigma_m,\sigma'_1,\dots,\sigma'_n)=\prod(1+t_i+t'_j) \]
	\end{thm}
	\begin{proof}
		The same as \cref{SW_tensor}. 
	\end{proof}

	\begin{defn}
		We define the dual bundle of $\omega$, $\Hom_\mathbb{C}(\omega,\mathbb{C})$ to be the bundle over the same space with each fibre identified with $\Hom_\mathbb{C}(F_x,\mathbb{C})$. 
	\end{defn}
	\begin{coro}
		If $\omega$ has a Hermitian metric, then $\Hom_\mathbb{C}(\omega,\mathbb{C})\simeq\bar{\omega}$. 
	\end{coro}
	\begin{proof}
		Follows from the vector space case. 
	\end{proof}

	We end this chapter again by calculating the Chern classes of the complex projective space. 
	\begin{prop}
		The total Chern class $c_{\tau_{\mathbb{CP}^n}}$ is equal to $(1-c_1(\gamma_1^1))^{n+1}$. Here, $\gamma_1^1$ is the canonical line bundle over $\mathbb{CP}^n$. 
	\end{prop}
	\begin{proof}
		Let $\gamma_1^1$ denote the canonical line bundle over $\mathbb{CP}^n$, and $\omega^n$ be its orthogonal complement in its tubular neighborhood. Again, we will show that the complex vector bundle $\Hom_{\mathbb{C}}(\gamma_1^1,\omega^n)\simeq\tau_{\mathbb{CP}^n}$. This is since a complex line $\ell$ through the origin of $\mathbb{C}^{n+1}$, the vector space $\Hom(\ell,\ell^\bot)$ can be identified with the neighborhood of $\ell$ in $\mathbb{CP}^n$ where $\ell'$ near $\ell$ is regarded as the graph of the function. \\
		Now adding the trivial bundle $\varepsilon^1$ to both sides of $\tau_{\mathbb{CP}^n}=\Hom(\gamma_1^1,\omega^1)$ we obtain 
		\[ \tau_{\mathbb{CP}^n}\oplus\varepsilon^1=\Hom(\gamma_1^1,\omega^1\oplus\gamma_1^1)=\Hom(\gamma_1^1,\varepsilon^1\oplus\cdots\oplus\varepsilon^1)=(\bar{\gamma_1^1})^{n+1} \]
		Hence we obtain the results immediately. 
	\end{proof}

\section{Chern-Weil theory}

	Interestingly, characteristic classes also arises in complex geometry, described in terms of curvature forms in the De-Rham cohomology. This characterization is often referred to as Chern-Weil theory. 

	\begin{defn}
		For a complex bundle $\zeta$, we define a connection on $\zeta$ is a $\mathbb{C}$-linear map 
		\[ \nabla:C^\infty(\zeta)\to C^\infty(\tau^*\otimes\zeta) \]
		satisfying the Leibniz formula $\nabla(fs)=\der f\otimes s+f\nabla(s)$. The image $\nabla(s)$ will be called the covariant derivative of $s$. 
	\end{defn}

	Intuitionally, the connection is very much like a derivative, and can be expressed locally. 
	\begin{thm}\label{Christoffel_symb}
		Suppose that $\zeta$ is a trivial $m$-bundle on a chart $U\subset M$, then we have $\nabla e_j=\omega^i_j\otimes e_i=\Gamma^i_{j,k}\der u^k\otimes e_i$. Here $\omega^i_j$ is called the connection forms. Since $\omega^i_j$ is not well-defined in coordinate changes, these 1-forms are only defined locally. \\
		Notice that the $m^2$ 1-forms $\omega^i_j$ form a (1,1)-tensor of 1-forms, hence is a matrix of 1-forms. We thus may regard $\omega$ as a section of $\tau^*_U\otimes\End(\zeta)$. 
	\end{thm}
	\begin{remark}
		The definition applies naturally to all principal $G$-bundles where $G$ is a Lie group. Notice that the case where $\zeta$ is a vector bundle is the special case where $G=\mathrm{GL}_n$, and $\End(\zeta)=\mathfrak{gl}(\zeta)$, the Lie algebra of $G$. This indicates that the connection of any principal $G$-bundle should be a section of $\tau^*_U\otimes\mathfrak{g}$. 
	\end{remark}

	\begin{defn}
		We define the curvature form of a connection $\nabla$ to be the matrix of 2-forms
		\[ \Omega^i_j=d\omega^i_j+\omega^i_k\wedge\omega^k_j \]
	\end{defn}
	\begin{prop}
		For tangent vector fields $X$ and $Y$, $\Omega(X,Y)=R(X,Y)\in\Gamma(\End(\zeta))$, where $R(X,Y)=\nabla_X\nabla_Y-\nabla_Y\nabla_X-\nabla_{[X,Y]}$ is the Riemann curvature tensor. 
	\end{prop}
	\begin{proof}
		We first calculate $\nabla_X\nabla_Y-\nabla_Y\nabla_X$. 
		\begin{align*}
			&(\nabla_X\nabla_Y - \nabla_Y\nabla_X)(a^ie_i) \\
			=&\nabla_X((Ya^j)e_j + a^k\langle Y,\omega^j_k\rangle e_j) - \nabla_Y((Xa^j)e_j + a^k\langle X,\omega^j_k\rangle e_j) \\
			=&(X(Ya^l) + (Ya^m)\langle X,\omega^l_m\rangle + X(a^k\langle Y,\omega^l_k\rangle) + a^k\langle Y,\omega^m_k\rangle\langle X,\omega^l_m\rangle)e_l \\
			-&(Y(Xa^l) + (Xa^m)\langle Y,\omega^l_m\rangle + Y(a^k\langle X,\omega^l_k\rangle) + a^k\langle X,\omega^m_k\rangle\langle Y,\omega^l_m\rangle)e_l \\
			=&(X(Ya^l) - Y(Xa^l))e_l \\
			+&(Y^j\pd{a^m}{u_j}X^k\Gamma^l_{m,k} - X^j\pd{a^m}{u_j}Y^k\Gamma^l_{m,k})e_l \\
			+&(X^j\pd{a^k}{u_j}Y^m\Gamma^l_{k,m} - Y^j\pd{a^k}{u_j}X^m\Gamma^l_{k,m})e_l \\
			+&(X^ja^k\pd{Y^m}{u_j}\Gamma^l_{k,m} - Y^ja^k\pd{X^m}{u_j}\Gamma^l_{k,m})e_l \\
			+&(X^ja^kY^m\pd{\Gamma^l_{k,m}}{u_j} - Y^ja^kX^m\pd{\Gamma^l_{k,m}}{u_j})e_l \\
			+&(a^kY^j\Gamma^m_{k,j}X^n\Gamma^l_{k,n} - a^kX^j\Gamma^m_{k,j}Y^n\Gamma^l_{k,n})
		\end{align*}
		Notice the second and third line cancels out. \\
		We then calculate $\nabla_{[X,Y]}$. 
		\begin{align*}
			&\nabla_{[X,Y]}(a^ie_i) \\
			=&([X,Y]a^l)e_l + a^k\langle[X,Y],\omega^l_k\rangle e_l \\
			=&(X(Ya^l) - Y(Xa^l))e_l + a^k\langle(X^j\pd{Y^m}{u_j}-Y^j\pd{X^m}{u_j})\pd{}{u_m},\Gamma^l_{k,m}\der u_m\rangle e_l \\
			=&(X(Ya^l) - Y(Xa^l))e_l + (X^ja^k\pd{Y^m}{u_j}\Gamma^l_{k,m} - Y^ja^k\pd{X^m}{u_j}\Gamma^l_{k,m})e_l
		\end{align*}
		which coincides the first line and the fourth line. \\
		We finally calculate $\Omega(X,Y)$. 
		\begin{align*}
			&\Omega(X,Y)(a^ie_i) \\
			=&(\der\omega^l_k(X,Y)+\omega^l_m\wedge\omega^m_k(X,Y))a^ke_l \\
			=&\pd{\Gamma^l_{k,m}}{u_j}(X^jY^m-Y^jX^m)e_l + (\Gamma^l_{m,n}X^n\Gamma^m_{k,j}Y^j - \Gamma^l_{m,n}Y^n\Gamma^m_{k,j}X^j) e_l
		\end{align*}
		which coincides with the fifth line and the sixth line.  
	\end{proof}
	\begin{coro}
		$\Omega$ is a globally defined matrix of 2-forms. 
	\end{coro}
	\begin{proof}
		$\nabla_X\nabla_Y-\nabla_Y\nabla_X-\nabla_{[X,Y]}$ is obviously globally defined. 
	\end{proof}

	Then, the curvature form satisfies the following Bianchi equality. 
	\begin{thm}\label{Bianchi}
		We have $\der\Omega^i_k=\Omega^i_j\wedge\omega^j_k-\omega^i_j\wedge\Omega^j_k$. We denote this by $\der\Omega=[\Omega,\omega]$. 
	\end{thm}
	\begin{proof}
		\begin{align*}
			&\der\Omega^i_k \\
			=&\der(\der\omega^i_k+\omega^i_j\wedge\omega^j_k) \\
			=&(\der\omega^i_j)\wedge\omega^j_k - \omega^i_j\wedge(\der\omega^j_k) \\
			=&(\Omega^i_j-\omega^i_l\wedge\omega^l_j)\wedge\omega^j_k-\omega^i_j\wedge(\Omega^j_k-\omega^j_l\wedge\omega^l_k) \\
			=&\Omega^i_j\wedge\omega^j_k - \omega^i_j\wedge\Omega^j_k \\
			+&\omega^i_j\wedge\omega^j_l\wedge\omega^l_k - \omega^i_l\wedge\omega^l_j\wedge\omega^j_k
		\end{align*}
		and the final line vanishes. 
	\end{proof}

	We then recall some basic definitions from linear algebra. 
	\begin{defn}
		An invariant polynomial of $n\times n$-matrix is a polynomial $P:\mathbb{C}^{n\times n}\simeq M_{n\times n}(\mathbb{C})\to\mathbb{C}$, satisfying the following equation for all matrices $A$ and invertible matrices $T$
		\[ P(TAT^{-1})=P(A) \]
	\end{defn}
	\begin{thm}\label{Struct_inv_poly}
		$P(A)$ is of the form $P(A)=p(\Tr A,\cdots,\Tr A^n)$. 
	\end{thm}
	\begin{proof}
		$P$ is a symmetric polynomial of the $n$ eigenvalues of $A$. 
	\end{proof}
	\begin{coro}
		The invariant polynomials are generated by the elementary symmetric polynomials of the eigenvalues, explicitly 
		\[ \sigma_k(A)=\sum_{1\leq i_1<\cdots<i_k\leq n}\lambda_{i_1}\cdots\lambda_{i_k} \]
	\end{coro}

	Notice for any invariant polynomial $P$, we have a map $1\otimes P:\Gamma(\Omega^2\otimes\End(\zeta))\to\Gamma(\Omega^2)$. 
	\begin{prop}
		$P(\Omega)\in\Omega^{2*}(M,\mathbb{C})$ is a closed form, where $\Omega$ is the curvature form. 
	\end{prop}
	\begin{proof}
		From \cref{Struct_inv_poly} we see that it suffices to prove $\der\Omega^k=0$. By \cref{Bianchi} and induction, we have $\der\Omega^k=[\Omega^k,\omega]$. Thus we have 
		\[ \der\Tr\Omega^k=\Tr\der\Omega^k=\Tr[\Omega^k,\omega]=0 \]
	\end{proof}
	\begin{prop}
		The cohomology class $[P(\Omega)]\in H^{2*}(M,\mathbb{C})$ does not depend on the choice of the connection. 
	\end{prop}
	\begin{proof}
		Let $\nabla^0$ and $\nabla^1$ be two connections on $\zeta$. We consider the vector bundle $\zeta'=p\times1:E\times\mathbb{R}\to M\times\mathbb{R}$, with connection $\nabla=t\tilde{\nabla}^1+(1-t)
		\tilde{\nabla}^0$, where $\tilde{\nabla}^i$ is defined by the pullback of the connection forms of $\nabla^i$ along $M\times\mathbb{R}\to M$. \\
		Let $i_0,i_1:M\to M\times\mathbb{R}$ be the inclusion at respectively $0$ and $1$. We let $\Omega^0$ and $\Omega^1$ be the curvature form of $\nabla^0$ and $\nabla^1$, then we have 
		\[ [P(\Omega^0)]=i_0^*[P(\Omega)]=i_1^*[P(\Omega)]=[P(\Omega^1)] \]
		where the middle equality is since $i_0$ and $i_1$ are homotopic. 
	\end{proof}

	Now, we have showed that for any invariant polynomial $P$, $[P(\Omega)]$ is a well-defined cohomology class. We will see that for certain invariant polynomial $P$, $[P(\Omega)]$ are exactly the Chern classes. Or more precisely, we have the following. 
	\begin{thm}\label{Chern_Weil}
		For a complex vector bundle $\zeta$ on a manifold $M$, we have 
		\[ c_k(\zeta)=\frac{1}{(2\pi i)^k}[\sigma_k(\Omega)] \]
	\end{thm}
	Before the theorem, we need the following lemma. 
	\begin{lemma}\label{Splitting_principle}
		Let $\zeta:E\to X$ be an $n$-dimensional complex vector bundle over a space $X$, then there exists a space equipped with a map $f:\tilde{X}\to X$ such that $f^*\zeta\simeq\ell_1\oplus\cdots\oplus\ell_n$, where $\ell_i$ are complex line bundles over $\tilde{X}$. Moreover, the map $f^*:H^*(X)\to H^*(\tilde{X})$ is injective. 
	\end{lemma}
	\begin{proof}
		We split off one line bundle at a time, and prove by induction. Consider the projective bundle $\mathbb{CP}^{n-1}\to P(\zeta)\xrightarrow{g}X$ corresponding to the complex vector bundle $\zeta$. Then $g^*\zeta$ has a canonical 1-dimensional subbundle $\ell_1$ with fibre over $(x,L)$ the line $L$. Then $g^*V=\ell_1\oplus\ell_1^\perp$. Then, notice $\mathbb{CP}^{n-1}$ has the cohomology ring a truncated polynomial ring, hence is free in all dimensions. Thus $g^*$ is injective by \cref{Leray-Hirsch}. We apply induction and immediately obtain the result. 
	\end{proof}
	\begin{proof}
		We now prove the theorem. We first prove the case for the first Chern class of a line bundle. Suppose $\zeta$ is a complex vector bundle over $M$, then suppose the classifying map of $\zeta$ is $\varphi:M\to\mathbb{CP}^\infty$. Since $E(\zeta)$ is a finite dimensional CW complex, by the cellular approximation theorem, we may homotope $\varphi$ to a map that factors through $\mathbb{CP}^N$. Moreover, we may assume $\varphi$ is smooth by some techniques in differential topology. Since $H^2(\mathbb{CP}^n;\mathbb{C})=\mathbb{C}c_1(\gamma_1^1)$. Then let $\Omega$ be the canonical connection of $\gamma_1^1$. We have $[\Omega]=\alpha c_1(\gamma_1^1)$. Then $[\Omega(\zeta)]=\alpha\varphi^*(c_1(\gamma_1^1))=\alpha c_1(\zeta)$. \\
		It then suffices to determin $\alpha$. We first focus on a 2-dimensional orientable Riemannian manifold $M$, and let $\tau_M$ be its tangent bundle. Then we have a canonical almost complex structure on $M$ given by rotating a vector in $\tang_x$ by $\pi/2$ counterclockwise. We denote the Levi-Civita connection on $\tau_M$ by $\nabla$. We then take a chart and let $e_1$ satisfy $|e_1|\equiv1$. Then let $e_2=ie_1$. Now, since length is preserved under parallel shifting, we have $\nabla_Xe_1\perp e_1$ for all $X$. Thus $\nabla e_1=\omega_1^2\otimes e_2$. Then, let the complexification of $\nabla$, $\nabla^\mathbb{C}$ be defined as $\nabla^\mathbb{C}s=\nabla s$ for all sections $s$. This is $\mathbb{C}$-linear since $\nabla$ is compatible with the Riemannian structure. Finally, from $\omega_1^2\otimes_\mathbb{C}e_2=i\omega_1^2\otimes e_1$ we see that the complex connection form $\omega^\mathbb{C}=i\omega_1^2$ (it is 1-dimensional). Thus the complex curvature form is $\Omega^\mathbb{C}=i\Omega_1^2$. \\
		We then apply the above calculations to the case $M=S^2$, equipped with the standard Riemann tensor. On one hand, by the pairing definition of integrals, we have
		\[ \int_{S^2}\Omega^\mathbb{C}=\langle\alpha c_1(\tau_{S^2}),\mu_{S^2}\rangle=\alpha\langle e(\tau_{S^2}),\mu_{S^2}\rangle=2\alpha \]
		On the other hand, by the above arguments, we have 
		\[ \int_U\Omega^\mathbb{C}=i\int_U\Omega_1^2=i\int_UK\der S=i\cdot\mathrm{Vol}(U) \]
		By taking the partition of unity, we have $2\alpha=i\cdot\mathrm{Vol}(S^2)=4\pi i$, hence $\alpha=2\pi i$. \\
		Finally, we treat the general case. Let $\zeta=p:E\to M$ be a $m$-dimensional complex vector bundle, and let 
		\[ P(A)=\sum_{k=1}^m\frac{1}{(2\pi i)^k}\sigma_k(A)=\det(I+\frac{A}{2\pi i}) \]
		By the theorem, we need to prove that this invariant polynomial corresponds to the total Chern class $c(\zeta)$. By \cref{Splitting_principle}, it suffices to prove the case where $\zeta=\ell_1\oplus\cdots\oplus\ell_m$. Then, we take a connection $\nabla^i$ on each line bundle $\ell_i$, and take the connection on $\zeta$ to be $\bigoplus\nabla^i$. In other words, we define the connection forms of $\nabla$ to be $\omega=\mathrm{diag}\{\omega^1,\dots,\omega^n\}$. Hence we have
		\begin{align*}
			&P(\Omega) \\
			=&\det\textrm{diag}\left(1+\frac{\Omega^1}{2\pi i},\dots,1+\frac{\Omega^m}{2\pi i}\right) \\
			=&P(\Omega^1)\cdots P(\Omega^m) \\
			=&c(\ell^1)\cdots c(\ell^m) \text{~by~the~case~for~line~bundles} \\
			=&c(E)
		\end{align*}
	\end{proof}
	As a corollary, we may define the Pontrajagin class using differential forms. 
	\begin{coro}
		For a vector bundle $\zeta=p:E\to M$ with a connection $\nabla$ and curvature $\Omega$, then the $k$-th Pontrajagin class can be characterized as follows. 
		\[ p_k(E)=\frac{1}{(2\pi)^2k}[\sigma_{2k}(\Omega)]\in H^{4k}(M;\mathbb{R}) \] 
	\end{coro}
	\begin{proof}
		Notice that a connection on $E$ induces a natural connection on $\mathbb{C}\otimes E$, with equal connection forms (one w.r.t $\mathbb{R}$ and another w.r.t $\mathbb{C}$). 
	\end{proof}

	As a final remark, we can prove the Chern-Gauss-Bonnet formula with this result. Before this, we first recall the Pfaffian of a skew-symmetric matrix. 
	\begin{defn}
		For an oriented $n$-dimensional inner product space, let $e_1,\dots,e_n$ be an orthonormal basis compatible with the orientation. Then we let $T$ be the linear map $T:\bigwedge^*V\to\mathbb{R}$ given by $T(\bigwedge^{<n}V)=0$ and $T(e_1\wedge\cdots\wedge e_n)=1$. \\
		For an element $A\in V\wedge V$, we may regard it as an element of $V\wedge V^*$ through the natural isomorphism given by the inner product, and hence is a skew-symmetric matrix. We define the Pfaffian of $A$ to be 
		\[ \mathrm{pf}(A)=T(\sum_{k=0}^\infty\frac{A^{\wedge k}}{k!}) \]
	\end{defn}
	\begin{prop}
		We have $(\mathrm{pf}(A))^2=\det A$. 
	\end{prop}
	\begin{proof}
		Directly check using the block diagonalization of $A$. 
	\end{proof}
	Finally we recall that the curvature form of a metric is skew-symmetric. Then we have the following. 
	\begin{thm}\label{Chern_Gauss_Bonnet}
		Let $\zeta=p:E\to M$ be a real vector bundle of dimension $2n$, then we have 
		\[ e(E)=[(\frac{-1}{2\pi})^n\mathrm{pf}(\Omega)]\in H^{2n}(M) \]
		In particular, we have
		\[ \chi(M)=\frac{-1}{(2\pi)^n}\int_M\mathrm{pf}(\Omega) \]
	\end{thm}
	\begin{proof}
		The second assertion follows immediately from the first one. Notice $e(\zeta)^2=p_n(\zeta)$, and $\mathrm{pf}(\Omega)^2=\det(\Omega)=\sigma_{2n}(\Omega)$. Thus we have 
		\[ e(\zeta)=\pm\left[\left(\frac{1}{2\pi}\right)^n\mathrm{pf}(\Omega)\right] \]
		Now it suffices to determin the sign. Notice that we may again pass to the canonical bundle $\eta^{2n}$ of $Gr_{2n}(\mathbb{R}^N)$ as in the proof of \cref{Chern_Weil}, so that the sign is independent of the choice of the manifold or vector bundles. \\
		When $n=1$, again by considering $S^2$ we have the sign a minus sign. For general $n$, we again use the splitting principle to reduce to the line bundle case. Therefore the sign is $(-1)^n$. 
	\end{proof}

\section{Pontrajagin classes}

	\begin{defn}
		For a real vector space $V$, we may obtain its complexification $V\otimes\mathbb{C}\simeq V\oplus iV$. Similarly, for a real vector bundle $\xi$ over some space $B$, we define its complexification to be the complex bundle $\xi\otimes\mathbb{C}\simeq\xi\oplus i\xi$ with the underlying real bundle $\xi\oplus\xi$ and the obvious complex structure. 
	\end{defn}

	\begin{lemma}
		For a real bundle $\xi$, we have $\xi\otimes\mathbb{C}\simeq\overline{\xi\otimes\mathbb{C}}$. 
	\end{lemma}
	\begin{proof}
		We directly map $\xi\otimes\mathbb{C}\to\overline{\xi\otimes\mathbb{C}}$ by $1_\xi\oplus-1_\xi$. This is an orientation preserving isomorphism. 
	\end{proof}
	
	We thus have the following condition. 
	\begin{coro}
		$c_{2i+1}(\xi\otimes\mathbb{C})$ is of order $2$. 
	\end{coro}
	\begin{proof}
		By \cref{Chern_ori} and the above lemma. 
	\end{proof}
	Therefore the following definition discarding the odd Chern classes makes sense. 
	\begin{defn}\label{Pontr_class}
		For a real bundle $\xi$ over $B$, we define its $i$-th Pontrajagin class $p_i(\xi)\in H^{4i}(B;\mathbb{Z})$ to be the class $(-1)^ic_{2i}(\xi\otimes\mathbb{C})$. Then the total Pontrajagin class is defined as $p(\xi)=1+p_1(\xi)+\cdots+p_{\lfloor n/2\rfloor}(\xi)$, where $\lfloor n/2\rfloor$ is the greatest integer smaller than $n/2$. 
	\end{defn}

	\begin{coro}\label{Pontr_nat_triv}
		Pontrajagin classes are natural w.r.t bundle maps. Moreover, $p(\xi\oplus\varepsilon^k)=p(\xi)$. 
	\end{coro}
	\begin{proof}
		Follows from \cref{Chern_nat} and \cref{Chern_prod}. 
	\end{proof}

	The product formula, however, does not hold for most general cases since we ignored the odd dimensional classes. We still have the weaker version below though. 
	\begin{coro}
		For real bundles $\xi$ and $\eta$, we have $2(p(\xi\oplus\eta)-p(\xi)p(\eta))=0$. 
	\end{coro}
	\begin{proof}
		Every term in the difference between them contains an odd ordered Chern class. 
	\end{proof}

	However, we do have the tensor product formula. 
	\begin{const}\label{Pontrajagin_tensor}
		By \cref{Chern_tensor}, we have $p_i(\xi\otimes\eta)=c_{2i}((\xi\otimes\eta)\otimes\mathbb{C})=c_{2i}((\xi\otimes\mathbb{C})\otimes(\eta\otimes\mathbb{C}))$. The rest follows from the tensor product formula. 
	\end{const}

	\begin{lemma}\label{Split_cplx}
		For a complex bundle $\omega$, we have $\omega_{\mathbb{R}}\otimes\mathbb{C}\simeq\omega\oplus\bar{\omega}$. 
	\end{lemma}
	\begin{proof}
		We consider a complex $n$-dimensional vector space $V$. Then this condition is equivalent to $V=W\oplus iW$ for some real $n$-dimensional vector space $W$ with complex structure $(x,y)\mapsto(-y,x)$. Therefore we have $V\oplus\bar{V}=W\oplus iW\oplus W\oplus -iW=V_{\mathbb{R}}\otimes\mathbb{C}$
	\end{proof}
	\begin{coro}
		For any $n$-dimensional complex vector bundle $\omega$ over $B$, the Chern classes $c_i(\omega)$ and the Pontrajagin classes $p_i(\omega_\mathbb{R})$ are related by the following formula
		\[ 1-p_1+\cdots+(-1)^np_n=(1-c_1+\cdots+(-1)^nc_n)(1+c_1+\cdots+c_n) \]
	\end{coro}
	\begin{proof}
		Follows from the preceding lemma, \cref{Chern_prod} and \cref{Chern_ori}. 
	\end{proof}
	
	\begin{lemma}
		The real $2n$-plane bundle $(\xi\otimes\mathbb{C})_\mathbb{R}$ is isomorphic to $\xi\oplus\xi$ under an isomorphism which either preserves or reverses orientation according to whether $n(n-1)/2$ is even or odd. 
	\end{lemma}
	\begin{proof}
		Direct by aligning the basis in the order $v_1,\dots,v_n,iv_1,\dots,iv_n$. 
	\end{proof}
	\begin{coro}
		If $\xi$ is an oriented $2k$-plane bundle, then the Pontrajagin class $p_k(\xi)$ is equal to the square of the Euler class $e(\xi)$. 
	\end{coro}
	\begin{proof}
		By definition, $p_k(\xi)=(-1)^kc_{2k}(\xi\otimes\mathbb{C})=(-1)^ke((\xi\otimes\mathbb{C})_\mathbb{R})$. But by the preceding lemma and \cref{Euler_prod}, we have $e(\xi\oplus\xi)=(-1)^{2k(2k-1)/2}e(\xi)^2=(-1)^ke(\xi)^2$. 
	\end{proof}

	We end this section with a simple application of the Pontrajagin classes. 
	\begin{thm}
		There are no almost complex structures on the sphere $S^{4k}$. 
	\end{thm}
	\begin{proof}
		The proof is by contradiction. If an almost complex structure $J$ exists on $S^{4k}$, then by \cref{Split_cplx} we have $\tau_{S^{4k}}\otimes\mathbb{C}\simeq\tau_{S^{4k}}\oplus\overline{\tau_{S^{4k}}}$, where the conjugation is given by the complex structure. 
		Notice that all cohomology except for the $0$-th and $4k$-th degree vanish, hence all Pontrajagin classes and Chern classes vanish for $\tau_{S^{4k}}$ except for the bottom and top term. Then we have
		\begin{align*}
			&1+(-1)^kp_k(\tau_{S^{4k}}) \\
			=&c(\tau_{S^{4k}}\otimes\mathbb{C}) \\
			=&c(\tau_{S^{4k}}\oplus\overline{\tau_{S^{4k}}}) \\
			=&c(\tau_{S^{4k}})c(\overline{\tau_{S^{4k}}}) \\
			=&(1+c_{2k}(\tau_{S^{4k}}))^2~\text{(by~\cref{Chern_ori})} \\
			=&1+2c_{2k}(\tau_{S^{4k}}) \\
			=&1+2e(\tau_{S^{4k}}). 
		\end{align*}
		Now, by \cref{Euler_class_char}, we have $p_k(\tau_{S^4k})$ non-trivial. However, since the normal bundle of the inclusion $S^n\inj\mathbb{R}^{n+1}$ is trivial, we have the Pontrajagin classes of $\tau_{S^{4k}}$ vanishes by \cref{Pontr_nat_triv}. 
	\end{proof}

\section{Oriented Grassmannians}

	Recall from \cref{Defn_Grassmannian} the definition of $Gr_n(\mathbb{R})$. We define in this section $\widetilde{Gr}_n(\mathbb{R})$, the oriented Grassmannian. 
	\begin{defn}
		For positive integers $n$ and $k$, we define $\widetilde{Gr}_n(\mathbb{R}^{n+k})$ to be the quotient space $\Hom_\mathbb{R}(\mathbb{R}^n,\mathbb{R}^{n+k})/R'$, where $R'$ is the equivalence relation where $f=g$ iff $\img f=\img g$ and $f$ and $g$ has the same orientation. We define $\widetilde{Gr}_n(\mathbb{R}):=\widetilde{Gr_n}(\mathbb{R}^\infty)=\ilim_k\widetilde{Gr}_n(\mathbb{R}^{n+k})$. Likewise we have a canonical $n$-bundle on $\widetilde{Gr}_n$, which we will denote by $\tilde{\eta}_n$. 
	\end{defn}
	\begin{remark}
		It is easy to see that we have a 2-sheeted covering space map $\widetilde{Gr}_n(\mathbb{R})\to Gr_n(\mathbb{R})$. Then by connectedness of $\widetilde{Gr}_n$ and that $\pi_1(Gr_n(\mathbb{R}))=\mathbb{Z}/2\mathbb{Z}$, we see that $\widetilde{Gr}_n(\mathbb{R})$ is the universal covering space of $Gr_n(\mathbb{R})$. This implies that $\widetilde{Gr}_n(\mathbb{R})$ is equivalent to $BSO(n)$ (it is a classical result in algebraic topology that the universal cover of $BO(n)$ is $BSO(n)$), hence classifies oriented $n$-bundles. 
	\end{remark}

	Now, the $\mathbb{Z}/2\mathbb{Z}$ coefficient cohomology of $\widetilde{Gr}_n(\mathbb{R})$ is easy to compute. 
	\begin{thm}\label{SW_gen_ori}
		$H^*(\widetilde{Gr}_n(\mathbb{R});\mathbb{Z}/2\mathbb{Z})$ is the free polynomial algebra generated by the Stiefel-Whitney classes $w_2(\tilde{\eta}^n),\dots,w_n(\tilde{\eta}^n)$. 
	\end{thm}
	\begin{proof}
		We first recall the Gysin sequence from \cref{Gysin_seq}. Then notice that for any 2-sheeted covering $\tilde{B}\to B$ is a $O(1)$-principal bundle over $B$, hence is associated to a line bundle $\xi$ over $B$. Therefore by the Gysin sequence we have a long exact sequence
		\[ \cdots\lto H^{i-1}(Gr_n;\mathbb{Z}/2\mathbb{Z})\lto H^i(Gr_n;\mathbb{Z}/2\mathbb{Z})\lto H^i(\tilde{Gr}_n;\mathbb{Z}/2\mathbb{Z})\lto\cdots \]
		where the map $H^{i-1}\to H^i$ is given by the cup product with the first Stiefel-Whitney class of the principal bundle $w_1(\xi)$, and the map $H^i\to H^i$ is given by the pullback along the projection. Notice that by \cref{SW_gen}, $H^1(Gr_n;\mathbb{Z}/2\mathbb{Z})=\mathbb{Z}/2\mathbb{Z}$, hence $w_1(\xi)$ is either $w_1(\eta^n)$ or $0$. But if $w_1(\xi)=0$, then we will have a short exact sequence at the zeroth cohomology
		\[ 0\lto H^0(Gr_n)\lto H^0(\tilde{Gr}_n)\lto H^0(Gr_n)\lto 0 \]
		which would imply that $\tilde{Gr}_n$ is not connected, a contradiction. 
		Therefore $w_1(\xi)=w_1(\eta^n)$, which implies the result immediately. 
	\end{proof}

	The computation of $\mathbb{Z}$ coefficient cohomology is much more difficult, so we will not present the computation in this note (maybe I will find some reference on Bockstein spectral sequences later). Instead, we calculate the $\Lambda$ coefficient cohomology, where $\Lambda$ is an integral domain containing $1/2$ (e.g. $\mathbb{Z}[1/2]$). 
	\begin{thm}
		If $\Lambda$ is an integral domain containing $1/2$, then the cohomology ring \\
		$H^*(\widetilde{Gr}_{2m+1}(\mathbb{R});\Lambda)$ is a polynomial ring generated by the Pontrajagin classes \\
		$p_1(\tilde{\eta}^{2m+1}),\dots,p_m(\tilde{\eta}^{2m+1})$. Similarly $H^*(\widetilde{Gr}_{2m}(\mathbb{R});\Lambda)$ is a polynomial ring generated by \\
		$p_1(\tilde{\eta}^{2m}),\dots,p_{m-1}(\tilde{\eta}^{2m})$ and $e(\tilde{\eta}^{2m})$. \\
		In other words, $H^*(\widetilde{Gr}_n(\mathbb{R});\Lambda)$ is generated by the classes $p_1,\dots,p_{\lfloor n/2\rfloor}$ and $e$, with the possible relation that $e^2=p_{n/2}$ when $n$ is even. 
	\end{thm}
	\begin{proof}
		The proof is by induction on $n$. When $n=1$, notice that $Gr_1(\mathbb{R})\simeq\mathbb{RP}^\infty$, hence $\widetilde{Gr}_1(\mathbb{R})\simeq S^\infty$, which is contractible. \\
		Similar to the complex case, we have $H^*(E_0(\eta^n))\simeq H^*(\widetilde{Gr}_{n-1})$ by sending $(f,x)\in E_0$ to $\img f\cap x^\perp$ preserving orientation. We have hence a long exact sequence
		\[ \cdots\lto H^i(\widetilde{Gr}_n)\lto H^{i+n}(\widetilde{Gr}_n)\lto H^{i+n}(\widetilde{Gr}_{n-1})\lto H^{i+1}(\widetilde{Gr}_{n+1}) \]
		where the first map is given by $\smile e$, and the second map $\lambda$ sends $p_k(\tilde{\eta}^n)$ to $p_k(\tilde{\eta}^{n-1})$, similar to the complex case. \\
		If $n$ is even, then by the induction hypothesis, $\lambda$ is surjective, hence we have short exact sequences 
		\[ 0\lto H^i(\widetilde{Gr}_n)\lto H^{i+n}(\widetilde{Gr}_n)\lto H^{i+n}(\widetilde{Gr}_{n-1})\lto 0 \]
		Since the cohomology ring of $\widetilde{Gr}_{n-1}$ is a polynomial ring generated by $p_1,\dots,p_{n/2-1}$, we have obviously that $\widetilde{Gr}_n$ is a polynomial ring generated by $p_1,\dots,p_{n/2-1}$ and $e$. \\
		If $n$ is odd, then $\smile e$ is the zero map, since $e$ is of torsion $2$ and $\Lambda$ is integral and contains $1/2$. Hence we have short exact sequences
		\[ 0\lto H^j(\widetilde{Gr}_{2m+1})\lto H^j(\widetilde{Gr}_{2m})\lto H^{j-2m}(\widetilde{Gr}_{2m+1})\lto 0 \]
		Hence we have $H^*(\widetilde{Gr}_{2m+1})$ a subring of $H^*(\widetilde{Gr}_{2m})$. If we denote the graded subring \\
		of $H^*(\widetilde{Gr}_{2m})$ generated by $p_1,\dots,p_m$ by $A^*$, then clearly by the definition of $\lambda$ we have $A^*\subseteq\img\lambda$. We must prove the equality holds. \\
		By induction hypothesis, every element of $H^j(\widetilde{Gr}_{2m})$ can be written as a sum $a+ea'$, where $a\in A^j$ and $a'\in A^{j-2m}$ (here $e$ is the Euler class of $\tilde{\eta}^{2m}$ with $e^2=p_m$). Hence we have $\rank H^j(\widetilde{Gr}_{2m})=\rank A^j+\rank A^{j-2m}$. But by the short exact sequence we have \\
		$\rank H^j(\widetilde{Gr}_{2m})=\rank H^j(\widetilde{Gr}_{2m+1})+\rank H^{j-2m}(\widetilde{Gr}_{2m+1})$. Since $A^*\subseteq\img\lambda$, we have \\
		$\rank A^j\leq \rank H^j(\widetilde{Gr}_{2m+1})$, but by the two equations above, the equality must holds, which implies $A^*=\img\lambda$. 
	\end{proof}

\section{Chern numbers and Pontrajagin numbers}

	In this final section, we introduce a somewhat more combinatorial part of the theory. 
	\begin{defn}\label{Chern_numbers}
		For $n$ a positive integer, we define $I=i_1,\dots,i_r$ as an unordered $r$-tuple a partition of $n$ iff $i_1+\cdots+i_r=n$. For two partitions $I=i_1,\dots,i_r$ and $J=j_1,\dots,j_s$ respectively of $n$ and $m$, we define $IJ=i_1,\dots,i_r,j_1,\dots,j_s$ is a partition of $m+n$. A refinement of $I$ is of the form $I_1\dots I_r$, where $I_r$ is a partition of $i_r$. 
	\end{defn}
	\begin{defn}
		For a compact $n$-dimensional complex manifold $K$, and $I$ a partition of $n$, we define the $I$-th Chern number $c_I[K]=c_{i_1}\cdots c_{i_r}[K]$ to be the integer $\langle c_{i_1}(\tau_K)\cdots c_{i_r}(\tau_K),\mu_K\rangle$, where both sides lies in the $2n$-th dimensional (co)homology.
	\end{defn}
	\begin{remark}
		Recall that $\mathbb{CP}^n$ has $c_i(\tau)=\binom{n+1}{i}a^i$, where $\langle a^n,\mu\rangle=1$. Therefore we have 
		\[ c_{i_1}\cdots c_{i_r}[\mathbb{CP}^n]=\binom{n+1}{i_1}\cdots\binom{n+1}{i_r} \]
	\end{remark}
	\begin{defn}\label{Pontrajagin_numbers}
		For a smooth compact $4n$-dimensional manifold $M$, and $I$ a partition of $4n$, we can likewise define the Pontrajagin numbers. 
	\end{defn}
	We can likewise define the Stiefel-Whitney numbers, though we will not focus on the definition. 
	\begin{defn}\label{SW_numbers}
		For a smooth manifold $M$, and $I$ a partition of $n$, we can also likewise define the Stiefel-Whitney numbers. The Stiefel-Whitney numbers lie in $\mathbb{Z}/2\mathbb{Z}$, unlike the Chern numbers and the Pontrajagin numbers. 
	\end{defn}

	We then introduce the perhaps most important object in combinatorics. 
	\begin{defn}
		We let $\mathcal{S}\subset\mathbb{Z}[t_1,\dots,t_n]$ be the subring consisting of all symmetric polynomials. Recall that $\mathcal{S}=\mathbb{Z}[\sigma_1,\dots,\sigma_n]$, where $\sigma_i$ is the $i$-th elementary symmetric polynomial. 
	\end{defn}
	\begin{lemma}
		If we endow $\mathcal{S}$ with the natural grading with $t_i$ having degree 1, then $\mathcal{S}^k$ is generated additively by $\sigma_{i_1}\cdots\sigma_{i_r}$ for $I=i_1,\dots,i_r$ a partition of $k$. 
	\end{lemma}
	\begin{proof}
		Obvious. 
	\end{proof}
	\begin{lemma}
		The above basis could be written explicitly as 
		\[ \sum t_{a_1}^{i_1}\cdots t_{a_r}^{i_r} \]
		for $I=i_1,\dots,i_r$ a partition of $k$, and the summation is taken over all $k$-tuples. 
	\end{lemma}
	\begin{proof}
		Obvious. 
	\end{proof}
	\begin{defn}
		For a partition $I$ of $k$, and $n\ge k$, then we define $s_I$ as the unique polynomial satisfying the following equation
		\[ s_I(\sigma_1,\dots,\sigma_n)=\sum t_{a_1}^{i_1}\cdots t_{a_r}^{i_r} \]
		Notice the definition of $s_I$ extends to the case where $n\le k$, in this case the $\sigma_i$ for $i\ge n$ are replaced by $0$. 
	\end{defn}

	\begin{lemma}\label{s_I_prod_formula}
		Notice that for $I$ a partition of $k$ and an $n$-dimensional complex bundle $\omega$ on $B$, we have $s_I(c(\omega)):=s_I(c_1,\dots,c_n)\in H^{2k}(B;\mathbb{Z})$. Then we have for $\omega$ and $\omega'$ on $B$
		\[ s_I(c(\omega\oplus\omega'))=\sum_{JK=I}s_J(c(\omega))s_K(c(\omega')) \]
	\end{lemma}
	\begin{proof}
		By definition, it suffices to prove that in $\mathbb{Z}[t_1,\dots,t_m,s_1,\dots,s_n]$ and a partition of $m+n$, $I$, we have (here $x_i$ is a substitution for $t_i$ and $s_j$)
		\[ \sum x_{a_1}^{i_1}\cdots x_{a_r}^{i_r}=\sum_{JK=I}\left(t_{b_1}^{j_1}\cdots t_{b_p}^{j_p}\right)\left(s_{c_1}^{k_1}\cdots s_{c_q}^{k_q}\right) \]
		Which is directly by definition. 
	\end{proof}
	\begin{coro}\label{s_k_additivity}
		$s_k(c(\omega\oplus\omega'))=s_k(c(\omega))+s_k(c(\omega'))$
	\end{coro}
	\begin{defn}
		According to the above defitions, for $I$ a partition of $n$, we define for an $n$-dimensional complex manifold $s_I(c)[K]$ to be the characteristic number $\langle s_I(c(\tau_K)),\mu_K\rangle$. 
	\end{defn}
	\begin{coro}\label{Chern_sym_poly}
		For complex manifolds $K$ and $L$, we have the following identity. 
		\[ s_I[K\times L]=\sum_{I_1I_2=I}s_{I_1}[K]s_{I_2}[L] \]
	\end{coro}
	\begin{proof}
		This is direct after noticing $\tau_{K\times L}=\pi_1^*(\tau_K)\oplus\pi_2^*(\tau_L)$. 
	\end{proof}
	\begin{coro}
		For any product $K\times L$ with $K$ $n$-dimensional and $L$ $m$-dimensional, we have $s_{m+n}(K\times L)=0$. 
	\end{coro}
	\begin{proof}
		$m+n$ admits no non-trivial decomposition into juxtapositions. 
	\end{proof}
	\begin{coro}
		$\mathbb{CP}^n$ is not a product of two manifolds. 
	\end{coro}
	\begin{proof}
		It is obvious that $s_n[\mathbb{CP}^n]=n+1\neq0$. 
	\end{proof}

	Finally, we prove that, in some sense, the Chern numbers are linearly independent. 
	\begin{thm}\label{Chern_num_indep}
		Let $K^i$ be an $i$-dimensional compact complex manifold with $s_i(c)[K^i]\neq0$. Then the $p(n)\times p(n)$ matrix ($p(n)$ is the number of partitions of $n$)
		\[ \left[ c_{i_1}\dots c_{i_r}[K^{j_1}\times\cdots\times K^{j_s}]\right] \]
		is non-singular. Here $I=i_1,\dots,i_r$ and $J=j_1,\dots,j_s$ runs through all partitions of $n$. 
	\end{thm}
	\begin{proof}
		By the definition of $s_I(\sigma_1,\dots,\sigma_n)$, we see that they form another basis of $\mathcal{S}^n$, hence it suffices to verify that the matrix
		\[ \left[ s_I(c)[K^{j_1}\times\cdots\times K^{j_s}]\right] \]
		is non-singular. As an immediate generalization of \cref{Chern_sym_poly}, we have
		\[ s_I(c)[K^{j_1}\times\cdots\times K^{j_s}]=\sum_{I_1\dots I_s=I}s_{I_1}[K^{j_1}]\cdots s_{I_s}[K^{j_s}] \]
		Note that for the product to be non-zero, $I_k$ must be a partition of $j_k$. Thus the number $s_I(c)[K^{j_1}\times\cdots\times K^{j_s}]$ is non-zero only if $I$ is a refinement of the partition $J=j_1,\dots,j_s$. Thus the matrix, under a suitable order, is triangular, with diagonal $s_{i_1,\dots,i_r}(c)[K^{i_1}\times\cdots\times K^{i_r}]=s_{i_1}(c)[K^{i_1}]\cdots s_{i_r}(c)[K^{i_r}]\neq 0$. Thus the matrix is non-singular. 
	\end{proof}
	We have an analogous theorem for Pontrajagin classes. 
	\begin{thm}\label{Pontrajagin_num_indep}
		Let $M^4,\dots,M^{4n}$ be oriented manifolds with $s_k(p)[M^{4k}]\neq 0$, then the $p(n)\times p(n)$ matrix 
		\[ p_{i_1}\cdots p_{i_r}[M^{4j_1}\times\cdots\times M^{4j_s}] \]
		is non-singular. 
	\end{thm}

	\begin{defn}
		For an $n$-dimensional complex vector bundle $\omega$ on $M$, we define its Chern character to be the formal sum
		\[ \Ch(\omega)=n+\sum_{k=1}^\infty\frac{s_k(c(\omega))}{k!}\in H^\Pi(M;\mathbb{Z}) \]
	\end{defn}
	\begin{thm}\label{Chern_char_additivity}
		We have $\Ch(\omega\oplus\omega')=\Ch(\omega)+\Ch(\omega')$. 
	\end{thm}
	\begin{proof}
		This is obvious by \cref{s_k_additivity}. 
	\end{proof}
	\begin{thm}\label{Chern_char_multiplicativity}
		We have $\Ch(\omega\otimes\omega')=\Ch(\omega)+\Ch(\omega')$. 
	\end{thm}
	\begin{proof}
		By \cref{Splitting_principle}, it suffices to consider the case where both $\omega$ and $\omega'$ are direct sums of line bundles. However, by the additivity above, it suffices to treat the case where both $\omega$ and $\omega'$ are line bundles. In this case, $c(\omega)=1+c_1(\omega)$, and $\Ch(\omega)=\exp(c_1(\omega))$. Thus we obviously have
		$\Ch(\omega)\Ch(\omega')=\exp(c_1(\omega))\exp(c_1(\omega'))=\exp(c_1(\omega)+c_1(\omega'))=\exp(c_1(\omega\otimes\omega'))$. 
	\end{proof}
	As in Chern-Weil theory, we have a characterization of the Chern character by the curvature tensor. 
	\begin{thm}\label{Chern_characteristic_curvature}
		For a complex vector bundle $\zeta$ on a manifold $M$ with connection $\nabla$ and curvature $\Omega$, we have
		\[ \Ch(\zeta)=[\Tr\exp(\frac{\Omega}{2\pi i})] \]
	\end{thm}
	\begin{proof}
		Recall from \cref{Chern_Weil} that we have
		\[ c_k(\zeta)=\frac{1}{(2\pi i)^k}[\sigma_k(\Omega)] \]
		This implies that the eigenvalues of $\Omega$ (as a 2-form) divided by $2\pi i$, say $\lambda_1,\dots,\lambda_n$, acts as the generators of $H^*(M;\mathbb{C})=\mathbb{C}\left[[\lambda_1],\dots,[\lambda_n]\right]$. 
		Then 
		\[ \Ch(\zeta)=\sum_{k=0}^\infty\frac{\lambda_1^k+\cdots+\lambda_n^k}{k!}=\sum_{k=0}^\infty\frac{1}{k!}\Tr\left(\frac{\Omega}{2\pi i}\right)^k=\Tr\exp(\frac{\Omega}{2\pi i}) \]
	\end{proof}
	
\chapter{A little bit more on Pontrajagin classes}

\section{Cobordism rings}

	In this section, we make use of the Stiefel-Whitney numbers and the Pontrajagin numbers to develop some elementary theory on Cobordism rings. In this section, all homology and cohomology of $BO$ and $MO$ are assumed to have $\mathbb{Z}/2\mathbb{Z}$ coefficients. We shall very frequently use notations of stable homotopy theory in this section, which readers can refer to in the second part of this note. 
	\begin{defn}
		We define $\mathfrak{N}_*$ to be the graded ring, with $\mathfrak{N}_m$ consisting of all unoriented cobordism class of $m$-dimensional compact manifolds ($M$ and $N$ are said to belong to the same cobordism class iff there is some $m+1$-dimensional manifold with boundary $P$ such that $\partial P=M\coprod N$), with addition the disjoint union and multiplication the cartesian product. \\
		We similarly define $\Omega_*$ to be the graded ring with $\Omega_m$ consisting of all oriented cobordism class of $m$-dimensional compact manifolds (that is, the set of all manifolds equipped with an orientation form, and cobordisms are asked to respect the orientation of the boundary of the manifold). 
	\end{defn}
	\begin{defn}
		We define $MO=\ilim MO(n)$, where $MO(n)$ is the Thom space of $\eta^n$, the universal bundle. We likewise define $MSO$ and $MU$. 
	\end{defn}

	We first prove the following theorem fundamental to cobordism theory. 
	\begin{thm}\label{MO_unoriented_cobordism}
		We have $\mathfrak{N}_n=\ilim_{k}\pi_{n+k}(MO(k))$
	\end{thm}
	\begin{remark}
		Notice that if we let $MO$ be the spectrum with the $k$-th space $MO(k)$, then the right hand side is the $n$-th homotopy group of $MO$, $\pi_*(MO)$. Also, recall that $MO$ is a ring spectra, hence $\pi_*(MO)$ is a graded ring. By construction, we would have $\mathfrak{N}_*\simeq\pi_*(MO)$ isomorphic as a ring. Notice also that the mod 2 homology $H_*(MO)$ is also a graded ring. 
	\end{remark}
	Before the proof, we need some simple facts in differential topology. 
	\begin{defn}
		For a smooth map $f:M\to N$, and a submanifold $X\subset N$. We say $f$ is transverse to $X$, denoted by $f\pitchfork X$, if for every $x\in f^{-1}(X)$, we have $\der f_x(\tang_xM)+\tang_{f(x)}X=\tang_{f(x)}N$. 
	\end{defn}
	\begin{thm}\label{transverse_pullback}
		Let $f:M\to N$ be a smooth map and $X$ be a submanifold of $N$. If $f\pitchfork X$, then $f^{-1}(X)$ is a submanifold of $M$. 
	\end{thm}
	\begin{thm}\label{generic_transversality}
		Let $f:M\to N$ be a smooth map and $X$ be a submanifold of $N$. Then there exists a smooth map $g:M\to N$, $g\pitchfork X$, and $f$ is homotopic to $g$. 
	\end{thm}
	\begin{proof}
		The proof will consist of three steps. The first one is to construct a well-defined map of abelian groups $\alpha:\mathfrak{N}_n\to\pi_n(MO)$. The second one is to construct $\beta:\pi_n(MO)\to\mathfrak{N}_n$. Finally we will verify that $\beta$ is indeed the inverse of $\alpha$. \\
		Now, given a cobordism class $[M]\in\mathfrak{N}_n$, let $M\in[M]$ be a representative. By the Whitney embedding theorem, we have an embedding $e:M\to\mathbb{R}^{n+k}$. Then, by \cref{Tube_Nbd}, we suppose there is a tubular neighborhood $T$ of $e(M)$. Now, obviously $\bar{T}/\partial\bar{T}\simeq\Th(\nu_M)$, where $\nu_M$ is the $k$-dimensional normal bundle of $e$. Then, let $S^{n+k}$ be the one-point compactification of the $\mathbb{R}^{n+k}$ above, so that we have a map $\eta:S^{n+k}\to\bar{T}/\partial\bar{T}\simeq\Th(\nu_M)$, often called the Thom collapse map. Therefore we already have an element $\eta\in\pi_{n+k}(\Th(\nu_M))$. We now consider the classifying map of $\nu_M$, $\varphi$. From the bundle map $\nu_M\to\eta_k$ we have a map $\Th(\nu_M)\to\Th(\eta_k)$, therefore $\varphi_*(\eta)$ is an element of $\pi_{n+k}(MO(k))\subseteq\pi_n(MO)$. \\
		Now, we prove $\alpha$ is independent of the choice of the cobordism representative $M$. It suffices to show that if $M$ is null-bordant, then $\alpha[M]$ is null-homotopic. Let $M=\partial W$. Then the inclusion $T\inj\mathbb{R}^{n+k}\inj S^{n+k}$ extends to a tubular neighborhood of $W$, $L$, to the disc $D^{n+k+1}$ (especially note here that the normal bundle $\nu_W$ is also of dimension $k$). Notice then that the two collapse maps agree, so that we have the following commutative square. 
		% https://q.uiver.app/#q=WzAsNCxbMCwwLCJTXntuK2t9Il0sWzAsMSwiRF57bitrKzF9Il0sWzEsMCwiXFxUaChcXG51X00pIl0sWzEsMSwiXFxUaChcXG51X1cpIl0sWzAsMl0sWzIsM10sWzAsMV0sWzEsM11d
		\[\begin{tikzcd}
			{S^{n+k}} & {\Th(\nu_M)} \\
			{D^{n+k+1}} & {\Th(\nu_W)}
			\arrow[from=1-1, to=1-2]
			\arrow[from=1-1, to=2-1]
			\arrow[from=1-2, to=2-2]
			\arrow[from=2-1, to=2-2]
		\end{tikzcd}\]
		Also, since the $\nu_W$ restricts to $\nu_M$ on $M$ (directly consider the tubular neighborhood, as its homotopic to the total space $\nu$ of the normal bundle), the classifying map of $\nu_M$ factors through the classifying map of $\nu_W$, say $\psi$, by the inclusion $\iota:M\inj W$. Therefore we have $\varphi_*=\psi_*\circ\iota_*:\nu_M\to\nu_W\to\eta^k$, and the map commutes at the level of Thom spaces. Therefore, by construction, the map $\varphi_*\circ\eta$ is null-homotopic through the inclusion into $D^{n+k+1}$. \\
		Then we verify that $\alpha[M\coprod N]=\alpha[M]+\alpha[N]$. We let $i:M\inj\mathbb{R}^{n+k}$ and $j:N\inj\mathbb{R}^{n+k}$ be two embeddings. Since both $M$ and $N$ are compact, we may suppose that their tubular neighborhoods lie in the distinct hemisphere of $S^{n+k}$. We notice that by definition we have $\Th(\nu_M)\vee\Th(\nu_N)\simeq\Th(\nu_{M\coprod N})$. Then by the construction above, the map $S^{n+k}\to\Th(\nu_{M\coprod N})$ factors as below. 
		\[ S^{n+k}\to S^{n+k}\vee S^{n+k}\to\Th(\nu_M)\vee\Th(\nu_N)\simeq\Th(\nu_{M\coprod N}) \]
		After composing with the map $\Th(\nu)\to MO(k)$, we obtain the result directly. \\
		We then construct $\beta$. Suppose we have some $[\varphi]\in\pi_n(MO)$, and a representative $\varphi:S^{n+k}\to MO(k)$. Since $S^{n+k}$ is compact, $\varphi$ factors through $\Th(\eta^k_{k+m})$, the Thom space of the canonical bundle of $Gr_k(\mathbb{R}^{k+m})$, since $Gr_k(\mathbb{R})=\ilim Gr_k(\mathbb{R}^{k+m})$. Now let $\xi:Gr_k(\mathbb{R}^{k+m})\to\Th(\eta^k_{k+m})$ denote the zero section of the bundle. We view $\xi(Gr_k(\mathbb{R}^{k+m}))$ as an embedded submanifold of $\Th(\eta^k_{k+m})$ of codimension $K$. By \cref{generic_transversality}, we have a map $f:S^{n+k}\to\Th(\eta^k_{k+m})$ such that $f$ is smooth, $f\simeq\varphi$, and $f\pitchfork\xi(Gr_k(\mathbb{R}^{k+m}))$. Thus let $X=f^{-1}(\xi(Gr_k(\mathbb{R}^{k+m})))$. By \cref{transverse_pullback}, $X$ is a submanifold of $S^{n+k}$ of codimension $k$, hence an $n$-dimensional manifold. Notice $X$ is compact since $\xi(Gr_k(\mathbb{R}^{k+m}))$ is closed. We define $X=\beta[\varphi]$. \\
		We then verify that $\beta$ is well-defined. Suppose we are given homotopic maps $\varphi\simeq\psi\in\pi_n(MO)$. By definition WLOG $\varphi$ and $\psi$ factors through $\Th(\eta^k_{k+m})$, as well as the homotopy between $\varphi$ and $\psi$, say $h$. Then we may find a smooth map $\mu:I\times S^{n+k}\to\Th(\eta^k_{k+m})$ homotopic to $h$ relative $\partial I\times S^{n+k}$, and transverse to $\xi(Gr_k(\mathbb{R}^{k+m}))$. By \cref{transverse_pullback} again, we have $\mu^{-1}(\xi(Gr_k(\mathbb{R}^{k+m})))$ a cobordism between the manifolds determined by $\varphi$ and $\psi$. \\
		Finally we verify that $\beta$ is an inverse of $\alpha$. We first consider $\beta\circ\alpha$. Let $M$ be a arbitrary $n$-dimensional compact manifold. Recall $\alpha[M]$ is given by $f\eta:S^{n+k}\to\Th(\nu_M)\to MO(k)$. Notice that by compactness again, the latter map factors through some $\Th(\eta^k_{k+m})$. By construction, the collapse map is constant on the tubular neighborhood $T$, and the second map sends only $M$ to the zero section in the construction of $\beta$. Thus $\beta\alpha[M]=[M]$. We then consider the inverse. Let $[\varphi]:S^{n+k}\to MO(k)$, and suppose the map factors through $\Th(\eta^k_{k+m})$ for some $m$, and is transverse to the zero section. Let $M$ be the result of the construction, so that $[M]=\beta[\varphi]$. Notice $\Th(\eta^k_{k+m})-\infty$ is a tubular neighborhood around the image of the zero section, which pulls back to a tubular neighborhood of $M$ in the $S^{n+k}$. This implies that we have the map $\varphi:S^{n+k}\to\Th(\eta^k_{k+m})$ factors through $\Th(\nu_M)$ by the collapse map. Therefore we have $\varphi=\alpha[M]=\alpha\beta[\varphi]$. 
	\end{proof}
	\begin{coro}\label{MSO_oriented_cobordism}
		We have $\Omega_n\simeq\pi_n(MSO)$. 
	\end{coro}
	\begin{proof}
		The proof is analogous. 
	\end{proof}

	With the help of \cref{MO_unoriented_cobordism}, we may begin to analyze the structure of the cobordism groups. Before this, we first recall from \cref{Steenrod_op} the definitions and properties of the Steenrod algebra. 
	\begin{defn}
		We define $\mathcal{A}$ to be the graded associative (not commutative) $\mathbb{Z}/2\mathbb{Z}$-algebra with grade $\mathcal{A}^r$ consisting of all stable cohomology operations on $H^*(-,\mathbb{Z}/2\mathbb{Z})$ of degree $r$, with addition given by addition and multiplication given by composition. \\
		By definition, the mod 2 cohomology of $X$ is $[X,\Sigma^nH\mathbb{Z}/2\mathbb{Z}]$, hence by Yoneda's theorem, the degree $r$ operations are precisely $[H\mathbb{Z}/2\mathbb{Z},\Sigma^rH\mathbb{Z}/2\mathbb{Z}]\simeq H^r(H\mathbb{Z}/2\mathbb{Z},\mathbb{Z}/2\mathbb{Z})$. 
	\end{defn}
	\begin{thm}
		$\mathcal{A}$ is isomorphic to the $\mathbb{Z}/2\mathbb{Z}$ algebra generated by $Sq^i$ taking quotients with the Adem relations (recall from \cref{Steenrod_op} the definition of Adem relations). 
	\end{thm}
	\begin{const}
		Notice that $\mathcal{A}$ can be given a coproduct structure $\mu:\mathcal{A}\to\mathcal{A}\otimes\mathcal{A}$, explicitly by 
		\[ Sq^k\mapsto\sum_{i+j=k}Sq^i\otimes Sq^j \]
		This, along with checking other basic properties, endows $\mathcal{A}$ the structure of a hopf algebra. \\
		Then, the dual algebra $\mathcal{A}_*=\Hom(\mathcal{A},\mathbb{Z}/2\mathbb{Z})$ inherits a natural commutative hopf algebra structure (whose algebra structure is inherited from the coalgebra structure of $\mathcal{A}$). 
	\end{const}
	\begin{thm}\label{Struct_A_*}
		We define $Sq^{I_r}=Sq^{2^{r-1}}Sq^{2^{r-2}}\cdots Sq^1$. Then, $\mathcal{A}_*\simeq\mathbb{Z}/2\mathbb{Z}[\xi_r]$, where $\xi_r$ is dual to $Sq^{I_r}$. Moreover, the coalgebra structure of $\mathcal{A}_*$ is given by 
		\[ \psi:\xi_r\mapsto\sum_{i+j=r}\xi^{2^j}_i\otimes\xi_j \]
	\end{thm}
	\begin{thm}\label{homology_Thom_isom}
		For an oriented $n$-bundle $\xi$ over $B$, we have $H_{*+n}(E,E_0)\simeq H_*(E)$. Moreover, the isomorphism is given by $\alpha\mapsto\alpha/u$, where $u$ is the fundamental class of the orientation. 
	\end{thm}
	\begin{proof}
		It suffices to prove for the case of the trivial bundle. The rest follows from a Mayer-Vietoris argument and the fact that $H_*(B)\simeq\ilim_{K}H_*(K)$, where $K$ runs through all compact subsets of $B$. However the trivial bundle case is easy by the Kunneth formula of homology. 
	\end{proof}

	Before we compute the homotopy groups of $MO$, we first compute its mod 2 homology. 
	\begin{thm}
		$H_*(MO)=\mathbb{Z}/2\mathbb{Z}[a_1,a_2,\dots]$. 
	\end{thm}
	\begin{proof}
		By \cref{homology_Thom_isom} above, we have a stablized version $H_*(MO)\simeq H_*(BO)$. Moreover, by construction, we find that the Thom isomorphism are compatible with the ring spectra structure, hence it is an isomorphism of graded rings. Then it suffices to compute $H_*(BO)$. Notice the multiplication of $H_*(BO)$ induces the comultiplication on $H^*(BO)\simeq\mathbb{Z}/2\mathbb{Z}[w_1,\dots]$ given by 
		\[ w_i\mapsto\sum_{j+k=i}w_j\otimes w_k \]
		Then $H_*(BO)$ is thus its dual, and the computations gives $H_*(BO)=\mathbb{Z}/2\mathbb{Z}[b_i]$, where $b_i$ is the image of $w_1^i\in H_i(BO(1))$ under the inclusion $H_i(BO(1))\to H_i(BO)$. Thus if we let $a_i$ be the preimage of $b_i$ under the isomorphism, we obtain the results. 
	\end{proof}
	\begin{remark}
		Since $b_i$ comes from $w_1^i$, and there is a slant product involved in the Thom isomorphism, we have $a_i$ coming from $w_1^{i+1}$, under the natural map $H_*(MO(1))\to H_*(MO)$ and the homotopy equivalence $BO(1)\simeq MO(1)$. 
	\end{remark}
	Now, since the cohomology of is an $\mathcal{A}$ module, the homology is naturally an $\mathcal{A}_*$ comodule (notice that in mod 2 coefficients, $H_*$ is the dual of $H^*$ since the homology is finite dimensional). We will see that the calculation of this comodule structure of certain spaces finishes the computation of $\pi_*(MO)$. 
	\begin{lemma}\label{comod_struct_BO1}
		Suppose the comodule structure of $H_*(BO(1))$ is given by $\gamma:H_*(BO(1))\to\mathcal{A}_*\otimes H_*(BO(1))$, and $\gamma(w_1^i)=a_{i,j}\otimes w_1^j$. Then we have
		\begin{equation*}
		a_{i,1}=
		\begin{cases}
			\xi_r & \text{if~}i=2^r \\
			0     & \text{otherwise}
		\end{cases}
		\end{equation*}
		and $a_{i,i}=1$. 
	\end{lemma}
	\begin{proof}
		By direct computation. 
	\end{proof}
	\begin{thm}
		We let $N_*=\mathbb{Z}/2\mathbb{Z}[u_i]$ for all $i>1$ and $i\neq 2^i-1$, with the degree of $u_i$ being $i$. Then we define a homomorphism $f:H_*(MO)\to N_*$ by 
		\begin{equation*}
		f(a_i)=
		\begin{cases}
			0   & \text{if~} i=2^r-1 \\
			u_i & \text{otherwise}
		\end{cases}
		\end{equation*}
		Then the composite $g=(1\otimes f)\circ\gamma:H_*(MO)\to \mathcal{A}_*\otimes H_*(MO)\to\mathcal{A}_*\otimes N_*$ is an isomorphism of $\mathcal{A}$-comodules and $\mathbb{Z}/2\mathbb{Z}$-algebras. 
	\end{thm}
	\begin{proof}
		$g$ is clearly a homomorphism of $\mathcal{A}$-comodules and $\mathbb{Z}/2\mathbb{Z}$-algebras. We prove it is an isomorphism. Notice that both the source and target of $g$ are polynomial algebras with one generator of degree $i$ for all $i$, by the structure of $\mathcal{A}_*$ stated in \cref{Struct_A_*}, hence it suffices to prove that $g$ sends generators to generators. Since $a_i=\iota(w_1^i)$, we may use this to compute $\gamma(a_i)=\iota(\gamma(w_1^{i+1}))$. Then, let $I$ be the ideal of $\mathcal{A}_*\otimes H_*(MO)$ generated by elements that are products of at least two generators. Then we have by \cref{comod_struct_BO1}
		\begin{equation*}
		\gamma(a_i)=\iota(\gamma(w_1^{i+1}))\equiv
		\begin{cases}
			\xi_r\otimes 1+1\otimes a_i & \text{if~}i=2^r-1 \\
			1\otimes a_i & \text{otherwise}
		\end{cases}
		(\mod~I)
		\end{equation*}
		This proves that the generators are sent to generators, since $\xi_r\otimes 1+1\otimes a_i$ is send to $\xi_r\otimes 1$ after applying $1\otimes f$. 
	\end{proof}
	\begin{const}
		Notice that the cohomology of spheres are concentrated in two degrees, hence the Steenrod operations on spheres are trivial. Therefore, when we consider the Hurewicz homomorphism $h_{MO}:\pi_*(MO)\to H_*(MO)$, we have $\gamma\circ h_{MO}(x)=1\otimes h_{MO}(x)$. Thus we consider the inclusion $N_*\simeq\mathbb{Z}/2\mathbb{Z}\otimes N*\inj\mathcal{A}_*\otimes N_*$. 
	\end{const}
	Finally we are ready to prove the theorem. 
	\begin{thm}\label{computation_cobordism_ring}
		The map $h_{MO}:\pi_*(MO)\to H_*(MO)$ is injective. Moreover, we have $g\circ h_{MO}$ an isomorphism onto $N_*$, as shown in the construction. 
	\end{thm}
	\begin{proof}
		Since each $\alpha\in H^n(MO)$ is represented by a map $MO\to\Sigma^nH\mathbb{Z}/2\mathbb{Z}$, we consider the space $\bigvee\Sigma^iH\mathbb{Z}/2\mathbb{Z}$, with each component corresponding to a generator of the cohomology of $MO$. Thus we have a natural map $f:MO\to\bigvee\Sigma^iH\mathbb{Z}/2\mathbb{Z}$. By definition, this map induces an isomorphism on mod 2 cohomology hence on homology. Then, consider the sequence
		\[ 0\lto\mathbb{Z}/2\mathbb{Z}\lto\mathbb{Z}/2^{k+1}\mathbb{Z}\lto\mathbb{Z}/2^k\mathbb{Z}\lto 0 \]
		and by the five lemma we have the map induce an isomorphism on $\mathbb{Z}/2^k\mathbb{Z}$ homology, and hence on the direct limit, $\mathbb{Z}[1/2]$ homology. Then, notice $\pi_*(MO)$ is of 2 torsion, hence $H_*(MO;\mathbb{Q})$ and $H_*(MO;\mathbb{Z}[1/p])$ for $p$ an odd prime are trivial. Thus by the sequence
		\[ 0\lto\mathbb{Z}\lto\mathbb{Q}\lto\bigoplus_p\mathbb{Z}[1/p]\lto0 \]
		we apply the five lemma again and find that $f$ induces an isomorphism on $\mathbb{Z}$ homology. Thus by the spectrum version of the Hurewicz theorem that $f$ must be a weak homotopy equivalence. \\
		Therefore, $f$ induces an isomorphism between $\pi_*(MO)$ and $\pi_*(\bigvee\Sigma^iH\mathbb{Z}/2\mathbb{Z})$. However by definition, $\pi_*(\bigvee\Sigma^iH\mathbb{Z}/2\mathbb{Z})\simeq N_*$. Therefore $\pi_*(MO)=N_*$, and $h_{MO}\simeq h_{\bigvee\Sigma^iH\mathbb{Z}/2\mathbb{Z}}$ is injective. Finally the following sequence of monomorphisms gives an injection from $\pi_*(MO)$ to $N_*\subset\mathcal{A}_*\otimes N_*$, and hence a isomorphism since on each degree they have the same finite dimension. 
		\[ \pi_*(MO)\to H_*(MO)\to\mathcal{A}_*\otimes H_*(MO)\to\mathcal{A}_*\otimes N_* \]
	\end{proof}

	\begin{thm}\label{cobordism_SW_num}
		Two manifolds are cobordant if and only if all their Stiefel-Whitney numbers are the same. In particular, a manifold is null-bordant if and only if its Stiefel-Whitney numbers vanish. 
	\end{thm}
	\begin{proof}
		Now, since the Hurewicz map $h_{MO}:\pi_*(MO)\to H_*(MO)$ is injective, it suffices to prove that for a manifold $M$, 
	\end{proof}




\section{The Hirzebruch signature theorem}

	\begin{defn}
		Let $\Lambda$ be a ring. Then we let $A^*=\bigoplus A^n$ denote a graded $\Lambda$-algebra. We define $A^\Pi$ to be the commutative ring of all formal sums $a_0+a_1+\cdots$. Then, we denote the multiplicative group of formal sums to be $A^u$. \\
		Then, Let $K_n(x_1,\dots,x_n)$ denote a sequence of polynomials with coefficients in $\Lambda$ such that $K_n(x_1,\dots,x_n)$ is homogeneous of degree $n$ if we assign $x_i$ to have degree $i$. We now define for $a\in A^u$, $K(a)\in A^u$ to be $1+K_1(a_1)+K_2(a_1,a_2)+\cdots$. Such $K$ is called a multiplicative sequence if we have $K(a)K(b)=K(ab)$. 
	\end{defn}

	\begin{lemma}\label{Char_mult_seq}
		Given a formal series $f(t)=1+\lambda_1t+\lambda_2t^2+\cdots$, there exists a unique multiplicative sequence $K$ such that $K(1+t)=f(t)$ (here $K(1+t)=1+K_1(t)+K_2(t,0)+\cdots$). 
	\end{lemma}
	\begin{proof}
		We first consider uniqueness. Consider the polynomial ring $\Lambda[t_1,\dots,t_n]=A^*$ with the natural grading. Then let $\sigma=(1+t_1)\cdots(1+t_n)\in A^u$. Thus $K(\sigma)=K(1+t_1)\cdots K(1+t_n)=f(t_1)\cdots f(t_n)$. Then taking the homogeneous part of degree $n$, we see that $K_n(\sigma_1,\dots,\sigma_n)$ is completely determined by the power series $f(t)$. Since $\sigma_i$ are algebraically independent, $K_n$ is unique. \\
		Then we prove existence. We let $K_n(\sigma_1,\dots,\sigma_n)=\sum\lambda_Is_I(\sigma_1,\dots,\sigma_n)$, the right hand side to be summed over all partitions of $n$ (here $\lambda_i=\lambda_{i_1}\cdots\lambda_{i_n}$). As in \cref{s_I_prod_formula}, we have
		\[ s_I(ab)=\sum_{JK=I}s_J(a)s_K(b) \]
		Thus we have
		\begin{align*}
			K(ab)=&\sum\lambda_I s_I(ab) \\
			=&\sum_I\lambda_I\sum_{JK=I}s_J(a)s_K(b) \\
			=&\sum_{JK=I}\lambda_Js_J(a)\lambda_Ks_K(b) \\
			=&K(a)K(b)
		\end{align*}
	\end{proof}

	\begin{defn}
		For $K$ a multiplicative sequence of rational coefficients and a manifold $M$ of dimension $4n$, we define its $K$-genus to be the rational number 
		\[ K[M]=\langle K_n(p_1,\dots,p_n),\mu_M\rangle \]
		Notice we have $K[M]$ a rational linear combination of the Pontrajagin numbers of $M$. \\
		For $M$ of dimension $m\neq 4n$, we define its $K$-genus to be $0$. 
	\end{defn}
	\begin{defn}
		For a compact oriented manifold $M$ of dimension $4n$, we define its signature to be the signature of the symmetric matrix given by $H^{2n}(M;\mathbb{Q})\times H^{2n}(M;\mathbb{Q})\to\mathbb{Q}$, $(a,b)\mapsto\langle a\smile b,\mu_M\rangle$. 
	\end{defn}
	\begin{thm}\label{Signature_thm}
		Let $L_k$ be the multiplicative sequence of polynomials belonging to the power series
		\[ \frac{\sqrt{t}}{\tanh\sqrt{t}}=1+\frac{1}{3}t-\frac{1}{45}t^2+\cdots+\frac{(-1)^{k-1}2^{2k}B_k}{(2k)!}t^{k}+\cdots \]
		where $B_k$ are the Bernoulli numbers. Then we have $\sigma(M)=L[M]$ for any compact oriented manifold $M$ of dimension $4n$. 
	\end{thm}
	\begin{proof}
		
	\end{proof}

